% sec:geo-raytracing — Hamiltonian Ray Tracing
\section{Hamiltonian Ray Tracing}
\label{sec:geo-raytracing}

Building a black-hole image means tracing one ray per pixel from the
observer backward through the curved geometry to its source --- the
accretion disc, the celestial sphere, or the shadow.  Each ray solves
the same eight-component ODE,
\cref{eq:geo-xdot,eq:geo-pdot-cov}, with only the initial
direction varying from pixel to pixel.  The cost of the image is
therefore set by two quantities: the expense of each integration step
and the number of steps each ray requires.  Both depend on how the
integrator interacts with the spacetime geometry --- a coupling that
this section makes explicit.
\coderef{Geodesic.Solver}{integrateGeodesic}

\paragraph{The per-step evaluation.}

Each evaluation of the right-hand side
$(\dot{x}^\mu, \dot{p}_\mu)$ at the current state
$(x^\mu, p_\mu, \lambda)$ proceeds in three stages:
%
\begin{enumerate}
\item \textbf{Inverse metric.}  Evaluate $\inversemetric$ at the
  current position.  For Kerr--Schild spacetimes, the inverse is
  $\eta^{\mu\nu} - f\,l^\mu l^\nu$
  (\cref{eq:kerr-ks-inverse}) --- the same function $f$ and null
  covector $l_\mu$ as the forward metric, with a single sign change.
  For a general spacetime without this algebraic structure, a
  $4 \times 4$ symmetric-matrix inversion is required at each step.
\item \textbf{Metric derivatives.}  Compute the four derivative
  matrices $\pd{\mu}\,g_{\alpha\beta}$ by passing the polymorphic
  metric function an AD dual-number type
  (\cref{ch:autodiff}).  The AD library's Jacobian
  evaluation extracts all $4 \times 10 = 40$ independent partial
  derivatives from the metric code.
\item \textbf{Hamilton's equations.}  Form the velocity
  $\dot{x}^\mu = \inversemetric\,p_\nu$ from step~1, then contract
  to obtain $\dot{p}_\mu = \tfrac{1}{2}\,
  (\pd{\mu}\,g_{\alpha\beta})\,\dot{x}^\alpha\,\dot{x}^\beta$ using
  the derivatives from step~2.  The velocity computed here is reused
  in the momentum contraction --- exactly as in the code, where
  \texttt{xdot} appears in both the return value and the
  \texttt{contractDg} helper.
\end{enumerate}
%
The inverse metric is evaluated once per RHS call and never
differentiated --- a property emphasised in \cref{sec:geo-formulations}
that keeps the per-step cost independent of the complexity of the
matrix inversion.
\implnote{The function \texttt{makeRHS} in
\codemodule{Geodesic.Solver} wraps \texttt{hamiltonianRHS} by
evaluating the inverse metric at the current position before each
call, keeping the per-step logic in a single closure.}

\paragraph{The adaptive integrator.}

The codebase uses Dormand--Prince RK45 \cite{DormandPrince1980}: an
embedded Runge--Kutta method that produces a fifth-order solution and
a fourth-order companion for local error estimation.  Each accepted
step evaluates the right-hand side six times (seven stages, but the
last stage of one step becomes the first of the next --- the
``first same as last'' property that saves one evaluation per step).

Step-size adaptation keeps the local error below prescribed
tolerances.  For each component $y_i$ of the eight-dimensional state
vector $(x^\mu, p_\mu)$, the scaled error is
%
\begin{equation}
  e_i = \frac{\abs{\delta y_i}}
    {\epsilon_a + \epsilon_r\,\abs{y_i}} \,,
  \label{eq:rk45-error}
\end{equation}
%
where $\delta y_i$ is the difference between the fifth- and
fourth-order solutions, $\epsilon_a$ is the absolute tolerance, and
$\epsilon_r$ the relative tolerance.  The mixed scaling ensures that
components near zero are governed by $\epsilon_a$ while large
components are governed by $\epsilon_r$, preventing any single
component from dominating the step-size decision.  The step is
accepted when $\max_i\,e_i \leq 1$.  After acceptance, the next
step size scales as
%
\begin{equation}
  h_{\text{new}} = h \cdot
    \min\!\bigl(s_{\max},\,
      \max\bigl(s_{\min},\,
        S\,(\max_i\,e_i)^{-1/5}\bigr)\bigr) \,,
  \label{eq:rk45-stepcontrol}
\end{equation}
%
so that halving the error produces a modest increase in step size,
while a large error triggers a proportionate reduction.  The safety
factor $S = 0.9$ guards against optimistic scaling, and the clamps
$s_{\min} = 0.2$, $s_{\max} = 5$ prevent extreme jumps.  The
exponent $-1/5$ arises because the embedded error estimate converges
as $\mathcal{O}(h^5)$; on a rejected step, the exponent switches
to $-1/4$ for more aggressive shrinkage, since the error estimate is
less reliable when the step has exceeded the asymptotic regime.
The default tolerances $\epsilon_a = \epsilon_r = 10^{-10}$ yield
Hamiltonian constraint violations $\abs{\mathcal{H}} \lesssim \epsilon_a$
for typical null geodesics through the Kerr geometry
(\cref{sec:geo-conservation}).
\xrefnote{The Runge--Kutta family, order conditions, and
step-size control theory are developed systematically in
\cref{ch:ode-integration}.  This section presents only the details
needed to understand the ray-tracing pipeline.}

A natural question is why the codebase does not use a symplectic
integrator --- St\"ormer--Verlet or leapfrog --- which would preserve
$\mathcal{H}$ to within bounded oscillations rather than allowing
monotonic drift.  The answer is step-size adaptivity.  Standard
symplectic methods are fixed-step: their conservation properties
depend on exactly preserving the symplectic two-form, which adaptive
step-size selection destroys.  Near the photon sphere, where the
effective potential has a saddle point and neighbouring trajectories
diverge exponentially, a fixed step small enough for the strong-field
region wastes orders of magnitude of computation in the weak-field
approach and escape phases.  Adaptive Dormand--Prince concentrates
work where the dynamics demand it.
\physnote{Symplectic integrators preserve a \emph{nearby} Hamiltonian
exactly (backward error analysis), so their energy error oscillates
rather than drifting.  For geodesic integration, the adaptive
step-size advantage outweighs this benefit.}

\paragraph{Step-size behaviour near the photon sphere.}

The photon sphere at $r = 3M$ (Schwarzschild) is an unstable
equilibrium for null circular orbits (\cref{sec:schw-orbits}).  A ray
with impact parameter $b$ close to the critical value
$b_c = 3\sqrt{3}\,M$ (\cref{eq:critical-impact}) passes near this
equilibrium, and small perturbations in the trajectory grow
exponentially.  The Lyapunov exponent of the unstable circular orbit
is $\gamma_L = 1/(3\sqrt{3}\,M)$ in coordinate time
\cite{Cardoso2009}: two initially neighbouring rays diverge as
$\sim\!\exp(\gamma_L\,t)$, forcing the integrator to shrink $h$ to
maintain accuracy.  Once the ray exits the strong-field zone, the step
size recovers --- a factor of ten or more variation over a single ray
trace is typical.  The adaptive integrator handles this automatically;
a fixed-step method would need the smallest step size throughout.

\begin{figure}[htbp]
  \centering
  %% Creator: Matplotlib, PGF backend
%%
%% To include the figure in your LaTeX document, write
%%   \input{<filename>.pgf}
%%
%% Make sure the required packages are loaded in your preamble
%%   \usepackage{pgf}
%%
%% Also ensure that all the required font packages are loaded; for instance,
%% the lmodern package is sometimes necessary when using math font.
%%   \usepackage{lmodern}
%%
%% Figures using additional raster images can only be included by \input if
%% they are in the same directory as the main LaTeX file. For loading figures
%% from other directories you can use the `import` package
%%   \usepackage{import}
%%
%% and then include the figures with
%%   \import{<path to file>}{<filename>.pgf}
%%
%% Matplotlib used the following preamble
%%   \def\mathdefault#1{#1}
%%   \everymath=\expandafter{\the\everymath\displaystyle}
%%   \IfFileExists{scrextend.sty}{
%%     \usepackage[fontsize=10.000000pt]{scrextend}
%%   }{
%%     \renewcommand{\normalsize}{\fontsize{10.000000}{12.000000}\selectfont}
%%     \normalsize
%%   }
%%   \usepackage{fontspec}
%%   \usepackage{unicode-math}
%%   \setmainfont{STIX Two Text}[Numbers=OldStyle, Ligatures=TeX]
%%   \setmathfont{STIX Two Math}
%%   \setmonofont{Source Code Pro}[Scale=MatchLowercase]
%%   \ifdefined\pdftexversion\else  % non-pdftex case.
%%     \usepackage{fontspec}
%%     \setmainfont{DejaVuSerif.ttf}[Path=\detokenize{/Users/gvn/.pyenv/versions/3.12.11/lib/python3.12/site-packages/matplotlib/mpl-data/fonts/ttf/}]
%%     \setsansfont{DejaVuSans.ttf}[Path=\detokenize{/Users/gvn/.pyenv/versions/3.12.11/lib/python3.12/site-packages/matplotlib/mpl-data/fonts/ttf/}]
%%     \setmonofont{DejaVuSansMono.ttf}[Path=\detokenize{/Users/gvn/.pyenv/versions/3.12.11/lib/python3.12/site-packages/matplotlib/mpl-data/fonts/ttf/}]
%%   \fi
%%   \makeatletter\@ifpackageloaded{underscore}{}{\usepackage[strings]{underscore}}\makeatother
%%
\begingroup%
\makeatletter%
\begin{pgfpicture}%
\pgfpathrectangle{\pgfpointorigin}{\pgfqpoint{3.834000in}{1.991700in}}%
\pgfusepath{use as bounding box, clip}%
\begin{pgfscope}%
\pgfsetbuttcap%
\pgfsetmiterjoin%
\definecolor{currentfill}{rgb}{1.000000,1.000000,1.000000}%
\pgfsetfillcolor{currentfill}%
\pgfsetlinewidth{0.000000pt}%
\definecolor{currentstroke}{rgb}{1.000000,1.000000,1.000000}%
\pgfsetstrokecolor{currentstroke}%
\pgfsetdash{}{0pt}%
\pgfpathmoveto{\pgfqpoint{0.000000in}{0.000000in}}%
\pgfpathlineto{\pgfqpoint{3.834000in}{0.000000in}}%
\pgfpathlineto{\pgfqpoint{3.834000in}{1.991700in}}%
\pgfpathlineto{\pgfqpoint{0.000000in}{1.991700in}}%
\pgfpathlineto{\pgfqpoint{0.000000in}{0.000000in}}%
\pgfpathclose%
\pgfusepath{fill}%
\end{pgfscope}%
\begin{pgfscope}%
\pgfsetbuttcap%
\pgfsetmiterjoin%
\definecolor{currentfill}{rgb}{1.000000,1.000000,1.000000}%
\pgfsetfillcolor{currentfill}%
\pgfsetlinewidth{0.000000pt}%
\definecolor{currentstroke}{rgb}{0.000000,0.000000,0.000000}%
\pgfsetstrokecolor{currentstroke}%
\pgfsetstrokeopacity{0.000000}%
\pgfsetdash{}{0pt}%
\pgfpathmoveto{\pgfqpoint{0.484000in}{0.352669in}}%
\pgfpathlineto{\pgfqpoint{3.814000in}{0.352669in}}%
\pgfpathlineto{\pgfqpoint{3.814000in}{1.971700in}}%
\pgfpathlineto{\pgfqpoint{0.484000in}{1.971700in}}%
\pgfpathlineto{\pgfqpoint{0.484000in}{0.352669in}}%
\pgfpathclose%
\pgfusepath{fill}%
\end{pgfscope}%
\begin{pgfscope}%
\pgfsetbuttcap%
\pgfsetroundjoin%
\definecolor{currentfill}{rgb}{0.000000,0.000000,0.000000}%
\pgfsetfillcolor{currentfill}%
\pgfsetlinewidth{0.501875pt}%
\definecolor{currentstroke}{rgb}{0.000000,0.000000,0.000000}%
\pgfsetstrokecolor{currentstroke}%
\pgfsetdash{}{0pt}%
\pgfsys@defobject{currentmarker}{\pgfqpoint{0.000000in}{0.000000in}}{\pgfqpoint{0.000000in}{0.041667in}}{%
\pgfpathmoveto{\pgfqpoint{0.000000in}{0.000000in}}%
\pgfpathlineto{\pgfqpoint{0.000000in}{0.041667in}}%
\pgfusepath{stroke,fill}%
}%
\begin{pgfscope}%
\pgfsys@transformshift{0.635265in}{0.352669in}%
\pgfsys@useobject{currentmarker}{}%
\end{pgfscope}%
\end{pgfscope}%
\begin{pgfscope}%
\definecolor{textcolor}{rgb}{0.000000,0.000000,0.000000}%
\pgfsetstrokecolor{textcolor}%
\pgfsetfillcolor{textcolor}%
\pgftext[x=0.635265in,y=0.304058in,,top]{\color{textcolor}{\sffamily\fontsize{8.000000}{9.600000}\selectfont\catcode`\^=\active\def^{\ifmmode\sp\else\^{}\fi}\catcode`\%=\active\def%{\%}0}}%
\end{pgfscope}%
\begin{pgfscope}%
\pgfsetbuttcap%
\pgfsetroundjoin%
\definecolor{currentfill}{rgb}{0.000000,0.000000,0.000000}%
\pgfsetfillcolor{currentfill}%
\pgfsetlinewidth{0.501875pt}%
\definecolor{currentstroke}{rgb}{0.000000,0.000000,0.000000}%
\pgfsetstrokecolor{currentstroke}%
\pgfsetdash{}{0pt}%
\pgfsys@defobject{currentmarker}{\pgfqpoint{0.000000in}{0.000000in}}{\pgfqpoint{0.000000in}{0.041667in}}{%
\pgfpathmoveto{\pgfqpoint{0.000000in}{0.000000in}}%
\pgfpathlineto{\pgfqpoint{0.000000in}{0.041667in}}%
\pgfusepath{stroke,fill}%
}%
\begin{pgfscope}%
\pgfsys@transformshift{1.037042in}{0.352669in}%
\pgfsys@useobject{currentmarker}{}%
\end{pgfscope}%
\end{pgfscope}%
\begin{pgfscope}%
\definecolor{textcolor}{rgb}{0.000000,0.000000,0.000000}%
\pgfsetstrokecolor{textcolor}%
\pgfsetfillcolor{textcolor}%
\pgftext[x=1.037042in,y=0.304058in,,top]{\color{textcolor}{\sffamily\fontsize{8.000000}{9.600000}\selectfont\catcode`\^=\active\def^{\ifmmode\sp\else\^{}\fi}\catcode`\%=\active\def%{\%}20}}%
\end{pgfscope}%
\begin{pgfscope}%
\pgfsetbuttcap%
\pgfsetroundjoin%
\definecolor{currentfill}{rgb}{0.000000,0.000000,0.000000}%
\pgfsetfillcolor{currentfill}%
\pgfsetlinewidth{0.501875pt}%
\definecolor{currentstroke}{rgb}{0.000000,0.000000,0.000000}%
\pgfsetstrokecolor{currentstroke}%
\pgfsetdash{}{0pt}%
\pgfsys@defobject{currentmarker}{\pgfqpoint{0.000000in}{0.000000in}}{\pgfqpoint{0.000000in}{0.041667in}}{%
\pgfpathmoveto{\pgfqpoint{0.000000in}{0.000000in}}%
\pgfpathlineto{\pgfqpoint{0.000000in}{0.041667in}}%
\pgfusepath{stroke,fill}%
}%
\begin{pgfscope}%
\pgfsys@transformshift{1.438820in}{0.352669in}%
\pgfsys@useobject{currentmarker}{}%
\end{pgfscope}%
\end{pgfscope}%
\begin{pgfscope}%
\definecolor{textcolor}{rgb}{0.000000,0.000000,0.000000}%
\pgfsetstrokecolor{textcolor}%
\pgfsetfillcolor{textcolor}%
\pgftext[x=1.438820in,y=0.304058in,,top]{\color{textcolor}{\sffamily\fontsize{8.000000}{9.600000}\selectfont\catcode`\^=\active\def^{\ifmmode\sp\else\^{}\fi}\catcode`\%=\active\def%{\%}40}}%
\end{pgfscope}%
\begin{pgfscope}%
\pgfsetbuttcap%
\pgfsetroundjoin%
\definecolor{currentfill}{rgb}{0.000000,0.000000,0.000000}%
\pgfsetfillcolor{currentfill}%
\pgfsetlinewidth{0.501875pt}%
\definecolor{currentstroke}{rgb}{0.000000,0.000000,0.000000}%
\pgfsetstrokecolor{currentstroke}%
\pgfsetdash{}{0pt}%
\pgfsys@defobject{currentmarker}{\pgfqpoint{0.000000in}{0.000000in}}{\pgfqpoint{0.000000in}{0.041667in}}{%
\pgfpathmoveto{\pgfqpoint{0.000000in}{0.000000in}}%
\pgfpathlineto{\pgfqpoint{0.000000in}{0.041667in}}%
\pgfusepath{stroke,fill}%
}%
\begin{pgfscope}%
\pgfsys@transformshift{1.840597in}{0.352669in}%
\pgfsys@useobject{currentmarker}{}%
\end{pgfscope}%
\end{pgfscope}%
\begin{pgfscope}%
\definecolor{textcolor}{rgb}{0.000000,0.000000,0.000000}%
\pgfsetstrokecolor{textcolor}%
\pgfsetfillcolor{textcolor}%
\pgftext[x=1.840597in,y=0.304058in,,top]{\color{textcolor}{\sffamily\fontsize{8.000000}{9.600000}\selectfont\catcode`\^=\active\def^{\ifmmode\sp\else\^{}\fi}\catcode`\%=\active\def%{\%}60}}%
\end{pgfscope}%
\begin{pgfscope}%
\pgfsetbuttcap%
\pgfsetroundjoin%
\definecolor{currentfill}{rgb}{0.000000,0.000000,0.000000}%
\pgfsetfillcolor{currentfill}%
\pgfsetlinewidth{0.501875pt}%
\definecolor{currentstroke}{rgb}{0.000000,0.000000,0.000000}%
\pgfsetstrokecolor{currentstroke}%
\pgfsetdash{}{0pt}%
\pgfsys@defobject{currentmarker}{\pgfqpoint{0.000000in}{0.000000in}}{\pgfqpoint{0.000000in}{0.041667in}}{%
\pgfpathmoveto{\pgfqpoint{0.000000in}{0.000000in}}%
\pgfpathlineto{\pgfqpoint{0.000000in}{0.041667in}}%
\pgfusepath{stroke,fill}%
}%
\begin{pgfscope}%
\pgfsys@transformshift{2.242374in}{0.352669in}%
\pgfsys@useobject{currentmarker}{}%
\end{pgfscope}%
\end{pgfscope}%
\begin{pgfscope}%
\definecolor{textcolor}{rgb}{0.000000,0.000000,0.000000}%
\pgfsetstrokecolor{textcolor}%
\pgfsetfillcolor{textcolor}%
\pgftext[x=2.242374in,y=0.304058in,,top]{\color{textcolor}{\sffamily\fontsize{8.000000}{9.600000}\selectfont\catcode`\^=\active\def^{\ifmmode\sp\else\^{}\fi}\catcode`\%=\active\def%{\%}80}}%
\end{pgfscope}%
\begin{pgfscope}%
\pgfsetbuttcap%
\pgfsetroundjoin%
\definecolor{currentfill}{rgb}{0.000000,0.000000,0.000000}%
\pgfsetfillcolor{currentfill}%
\pgfsetlinewidth{0.501875pt}%
\definecolor{currentstroke}{rgb}{0.000000,0.000000,0.000000}%
\pgfsetstrokecolor{currentstroke}%
\pgfsetdash{}{0pt}%
\pgfsys@defobject{currentmarker}{\pgfqpoint{0.000000in}{0.000000in}}{\pgfqpoint{0.000000in}{0.041667in}}{%
\pgfpathmoveto{\pgfqpoint{0.000000in}{0.000000in}}%
\pgfpathlineto{\pgfqpoint{0.000000in}{0.041667in}}%
\pgfusepath{stroke,fill}%
}%
\begin{pgfscope}%
\pgfsys@transformshift{2.644151in}{0.352669in}%
\pgfsys@useobject{currentmarker}{}%
\end{pgfscope}%
\end{pgfscope}%
\begin{pgfscope}%
\definecolor{textcolor}{rgb}{0.000000,0.000000,0.000000}%
\pgfsetstrokecolor{textcolor}%
\pgfsetfillcolor{textcolor}%
\pgftext[x=2.644151in,y=0.304058in,,top]{\color{textcolor}{\sffamily\fontsize{8.000000}{9.600000}\selectfont\catcode`\^=\active\def^{\ifmmode\sp\else\^{}\fi}\catcode`\%=\active\def%{\%}100}}%
\end{pgfscope}%
\begin{pgfscope}%
\pgfsetbuttcap%
\pgfsetroundjoin%
\definecolor{currentfill}{rgb}{0.000000,0.000000,0.000000}%
\pgfsetfillcolor{currentfill}%
\pgfsetlinewidth{0.501875pt}%
\definecolor{currentstroke}{rgb}{0.000000,0.000000,0.000000}%
\pgfsetstrokecolor{currentstroke}%
\pgfsetdash{}{0pt}%
\pgfsys@defobject{currentmarker}{\pgfqpoint{0.000000in}{0.000000in}}{\pgfqpoint{0.000000in}{0.041667in}}{%
\pgfpathmoveto{\pgfqpoint{0.000000in}{0.000000in}}%
\pgfpathlineto{\pgfqpoint{0.000000in}{0.041667in}}%
\pgfusepath{stroke,fill}%
}%
\begin{pgfscope}%
\pgfsys@transformshift{3.045928in}{0.352669in}%
\pgfsys@useobject{currentmarker}{}%
\end{pgfscope}%
\end{pgfscope}%
\begin{pgfscope}%
\definecolor{textcolor}{rgb}{0.000000,0.000000,0.000000}%
\pgfsetstrokecolor{textcolor}%
\pgfsetfillcolor{textcolor}%
\pgftext[x=3.045928in,y=0.304058in,,top]{\color{textcolor}{\sffamily\fontsize{8.000000}{9.600000}\selectfont\catcode`\^=\active\def^{\ifmmode\sp\else\^{}\fi}\catcode`\%=\active\def%{\%}120}}%
\end{pgfscope}%
\begin{pgfscope}%
\pgfsetbuttcap%
\pgfsetroundjoin%
\definecolor{currentfill}{rgb}{0.000000,0.000000,0.000000}%
\pgfsetfillcolor{currentfill}%
\pgfsetlinewidth{0.501875pt}%
\definecolor{currentstroke}{rgb}{0.000000,0.000000,0.000000}%
\pgfsetstrokecolor{currentstroke}%
\pgfsetdash{}{0pt}%
\pgfsys@defobject{currentmarker}{\pgfqpoint{0.000000in}{0.000000in}}{\pgfqpoint{0.000000in}{0.041667in}}{%
\pgfpathmoveto{\pgfqpoint{0.000000in}{0.000000in}}%
\pgfpathlineto{\pgfqpoint{0.000000in}{0.041667in}}%
\pgfusepath{stroke,fill}%
}%
\begin{pgfscope}%
\pgfsys@transformshift{3.447705in}{0.352669in}%
\pgfsys@useobject{currentmarker}{}%
\end{pgfscope}%
\end{pgfscope}%
\begin{pgfscope}%
\definecolor{textcolor}{rgb}{0.000000,0.000000,0.000000}%
\pgfsetstrokecolor{textcolor}%
\pgfsetfillcolor{textcolor}%
\pgftext[x=3.447705in,y=0.304058in,,top]{\color{textcolor}{\sffamily\fontsize{8.000000}{9.600000}\selectfont\catcode`\^=\active\def^{\ifmmode\sp\else\^{}\fi}\catcode`\%=\active\def%{\%}140}}%
\end{pgfscope}%
\begin{pgfscope}%
\pgfsetbuttcap%
\pgfsetroundjoin%
\definecolor{currentfill}{rgb}{0.000000,0.000000,0.000000}%
\pgfsetfillcolor{currentfill}%
\pgfsetlinewidth{0.301125pt}%
\definecolor{currentstroke}{rgb}{0.000000,0.000000,0.000000}%
\pgfsetstrokecolor{currentstroke}%
\pgfsetdash{}{0pt}%
\pgfsys@defobject{currentmarker}{\pgfqpoint{0.000000in}{0.000000in}}{\pgfqpoint{0.000000in}{0.027778in}}{%
\pgfpathmoveto{\pgfqpoint{0.000000in}{0.000000in}}%
\pgfpathlineto{\pgfqpoint{0.000000in}{0.027778in}}%
\pgfusepath{stroke,fill}%
}%
\begin{pgfscope}%
\pgfsys@transformshift{0.534821in}{0.352669in}%
\pgfsys@useobject{currentmarker}{}%
\end{pgfscope}%
\end{pgfscope}%
\begin{pgfscope}%
\pgfsetbuttcap%
\pgfsetroundjoin%
\definecolor{currentfill}{rgb}{0.000000,0.000000,0.000000}%
\pgfsetfillcolor{currentfill}%
\pgfsetlinewidth{0.301125pt}%
\definecolor{currentstroke}{rgb}{0.000000,0.000000,0.000000}%
\pgfsetstrokecolor{currentstroke}%
\pgfsetdash{}{0pt}%
\pgfsys@defobject{currentmarker}{\pgfqpoint{0.000000in}{0.000000in}}{\pgfqpoint{0.000000in}{0.027778in}}{%
\pgfpathmoveto{\pgfqpoint{0.000000in}{0.000000in}}%
\pgfpathlineto{\pgfqpoint{0.000000in}{0.027778in}}%
\pgfusepath{stroke,fill}%
}%
\begin{pgfscope}%
\pgfsys@transformshift{0.735710in}{0.352669in}%
\pgfsys@useobject{currentmarker}{}%
\end{pgfscope}%
\end{pgfscope}%
\begin{pgfscope}%
\pgfsetbuttcap%
\pgfsetroundjoin%
\definecolor{currentfill}{rgb}{0.000000,0.000000,0.000000}%
\pgfsetfillcolor{currentfill}%
\pgfsetlinewidth{0.301125pt}%
\definecolor{currentstroke}{rgb}{0.000000,0.000000,0.000000}%
\pgfsetstrokecolor{currentstroke}%
\pgfsetdash{}{0pt}%
\pgfsys@defobject{currentmarker}{\pgfqpoint{0.000000in}{0.000000in}}{\pgfqpoint{0.000000in}{0.027778in}}{%
\pgfpathmoveto{\pgfqpoint{0.000000in}{0.000000in}}%
\pgfpathlineto{\pgfqpoint{0.000000in}{0.027778in}}%
\pgfusepath{stroke,fill}%
}%
\begin{pgfscope}%
\pgfsys@transformshift{0.836154in}{0.352669in}%
\pgfsys@useobject{currentmarker}{}%
\end{pgfscope}%
\end{pgfscope}%
\begin{pgfscope}%
\pgfsetbuttcap%
\pgfsetroundjoin%
\definecolor{currentfill}{rgb}{0.000000,0.000000,0.000000}%
\pgfsetfillcolor{currentfill}%
\pgfsetlinewidth{0.301125pt}%
\definecolor{currentstroke}{rgb}{0.000000,0.000000,0.000000}%
\pgfsetstrokecolor{currentstroke}%
\pgfsetdash{}{0pt}%
\pgfsys@defobject{currentmarker}{\pgfqpoint{0.000000in}{0.000000in}}{\pgfqpoint{0.000000in}{0.027778in}}{%
\pgfpathmoveto{\pgfqpoint{0.000000in}{0.000000in}}%
\pgfpathlineto{\pgfqpoint{0.000000in}{0.027778in}}%
\pgfusepath{stroke,fill}%
}%
\begin{pgfscope}%
\pgfsys@transformshift{0.936598in}{0.352669in}%
\pgfsys@useobject{currentmarker}{}%
\end{pgfscope}%
\end{pgfscope}%
\begin{pgfscope}%
\pgfsetbuttcap%
\pgfsetroundjoin%
\definecolor{currentfill}{rgb}{0.000000,0.000000,0.000000}%
\pgfsetfillcolor{currentfill}%
\pgfsetlinewidth{0.301125pt}%
\definecolor{currentstroke}{rgb}{0.000000,0.000000,0.000000}%
\pgfsetstrokecolor{currentstroke}%
\pgfsetdash{}{0pt}%
\pgfsys@defobject{currentmarker}{\pgfqpoint{0.000000in}{0.000000in}}{\pgfqpoint{0.000000in}{0.027778in}}{%
\pgfpathmoveto{\pgfqpoint{0.000000in}{0.000000in}}%
\pgfpathlineto{\pgfqpoint{0.000000in}{0.027778in}}%
\pgfusepath{stroke,fill}%
}%
\begin{pgfscope}%
\pgfsys@transformshift{1.137487in}{0.352669in}%
\pgfsys@useobject{currentmarker}{}%
\end{pgfscope}%
\end{pgfscope}%
\begin{pgfscope}%
\pgfsetbuttcap%
\pgfsetroundjoin%
\definecolor{currentfill}{rgb}{0.000000,0.000000,0.000000}%
\pgfsetfillcolor{currentfill}%
\pgfsetlinewidth{0.301125pt}%
\definecolor{currentstroke}{rgb}{0.000000,0.000000,0.000000}%
\pgfsetstrokecolor{currentstroke}%
\pgfsetdash{}{0pt}%
\pgfsys@defobject{currentmarker}{\pgfqpoint{0.000000in}{0.000000in}}{\pgfqpoint{0.000000in}{0.027778in}}{%
\pgfpathmoveto{\pgfqpoint{0.000000in}{0.000000in}}%
\pgfpathlineto{\pgfqpoint{0.000000in}{0.027778in}}%
\pgfusepath{stroke,fill}%
}%
\begin{pgfscope}%
\pgfsys@transformshift{1.237931in}{0.352669in}%
\pgfsys@useobject{currentmarker}{}%
\end{pgfscope}%
\end{pgfscope}%
\begin{pgfscope}%
\pgfsetbuttcap%
\pgfsetroundjoin%
\definecolor{currentfill}{rgb}{0.000000,0.000000,0.000000}%
\pgfsetfillcolor{currentfill}%
\pgfsetlinewidth{0.301125pt}%
\definecolor{currentstroke}{rgb}{0.000000,0.000000,0.000000}%
\pgfsetstrokecolor{currentstroke}%
\pgfsetdash{}{0pt}%
\pgfsys@defobject{currentmarker}{\pgfqpoint{0.000000in}{0.000000in}}{\pgfqpoint{0.000000in}{0.027778in}}{%
\pgfpathmoveto{\pgfqpoint{0.000000in}{0.000000in}}%
\pgfpathlineto{\pgfqpoint{0.000000in}{0.027778in}}%
\pgfusepath{stroke,fill}%
}%
\begin{pgfscope}%
\pgfsys@transformshift{1.338375in}{0.352669in}%
\pgfsys@useobject{currentmarker}{}%
\end{pgfscope}%
\end{pgfscope}%
\begin{pgfscope}%
\pgfsetbuttcap%
\pgfsetroundjoin%
\definecolor{currentfill}{rgb}{0.000000,0.000000,0.000000}%
\pgfsetfillcolor{currentfill}%
\pgfsetlinewidth{0.301125pt}%
\definecolor{currentstroke}{rgb}{0.000000,0.000000,0.000000}%
\pgfsetstrokecolor{currentstroke}%
\pgfsetdash{}{0pt}%
\pgfsys@defobject{currentmarker}{\pgfqpoint{0.000000in}{0.000000in}}{\pgfqpoint{0.000000in}{0.027778in}}{%
\pgfpathmoveto{\pgfqpoint{0.000000in}{0.000000in}}%
\pgfpathlineto{\pgfqpoint{0.000000in}{0.027778in}}%
\pgfusepath{stroke,fill}%
}%
\begin{pgfscope}%
\pgfsys@transformshift{1.539264in}{0.352669in}%
\pgfsys@useobject{currentmarker}{}%
\end{pgfscope}%
\end{pgfscope}%
\begin{pgfscope}%
\pgfsetbuttcap%
\pgfsetroundjoin%
\definecolor{currentfill}{rgb}{0.000000,0.000000,0.000000}%
\pgfsetfillcolor{currentfill}%
\pgfsetlinewidth{0.301125pt}%
\definecolor{currentstroke}{rgb}{0.000000,0.000000,0.000000}%
\pgfsetstrokecolor{currentstroke}%
\pgfsetdash{}{0pt}%
\pgfsys@defobject{currentmarker}{\pgfqpoint{0.000000in}{0.000000in}}{\pgfqpoint{0.000000in}{0.027778in}}{%
\pgfpathmoveto{\pgfqpoint{0.000000in}{0.000000in}}%
\pgfpathlineto{\pgfqpoint{0.000000in}{0.027778in}}%
\pgfusepath{stroke,fill}%
}%
\begin{pgfscope}%
\pgfsys@transformshift{1.639708in}{0.352669in}%
\pgfsys@useobject{currentmarker}{}%
\end{pgfscope}%
\end{pgfscope}%
\begin{pgfscope}%
\pgfsetbuttcap%
\pgfsetroundjoin%
\definecolor{currentfill}{rgb}{0.000000,0.000000,0.000000}%
\pgfsetfillcolor{currentfill}%
\pgfsetlinewidth{0.301125pt}%
\definecolor{currentstroke}{rgb}{0.000000,0.000000,0.000000}%
\pgfsetstrokecolor{currentstroke}%
\pgfsetdash{}{0pt}%
\pgfsys@defobject{currentmarker}{\pgfqpoint{0.000000in}{0.000000in}}{\pgfqpoint{0.000000in}{0.027778in}}{%
\pgfpathmoveto{\pgfqpoint{0.000000in}{0.000000in}}%
\pgfpathlineto{\pgfqpoint{0.000000in}{0.027778in}}%
\pgfusepath{stroke,fill}%
}%
\begin{pgfscope}%
\pgfsys@transformshift{1.740152in}{0.352669in}%
\pgfsys@useobject{currentmarker}{}%
\end{pgfscope}%
\end{pgfscope}%
\begin{pgfscope}%
\pgfsetbuttcap%
\pgfsetroundjoin%
\definecolor{currentfill}{rgb}{0.000000,0.000000,0.000000}%
\pgfsetfillcolor{currentfill}%
\pgfsetlinewidth{0.301125pt}%
\definecolor{currentstroke}{rgb}{0.000000,0.000000,0.000000}%
\pgfsetstrokecolor{currentstroke}%
\pgfsetdash{}{0pt}%
\pgfsys@defobject{currentmarker}{\pgfqpoint{0.000000in}{0.000000in}}{\pgfqpoint{0.000000in}{0.027778in}}{%
\pgfpathmoveto{\pgfqpoint{0.000000in}{0.000000in}}%
\pgfpathlineto{\pgfqpoint{0.000000in}{0.027778in}}%
\pgfusepath{stroke,fill}%
}%
\begin{pgfscope}%
\pgfsys@transformshift{1.941041in}{0.352669in}%
\pgfsys@useobject{currentmarker}{}%
\end{pgfscope}%
\end{pgfscope}%
\begin{pgfscope}%
\pgfsetbuttcap%
\pgfsetroundjoin%
\definecolor{currentfill}{rgb}{0.000000,0.000000,0.000000}%
\pgfsetfillcolor{currentfill}%
\pgfsetlinewidth{0.301125pt}%
\definecolor{currentstroke}{rgb}{0.000000,0.000000,0.000000}%
\pgfsetstrokecolor{currentstroke}%
\pgfsetdash{}{0pt}%
\pgfsys@defobject{currentmarker}{\pgfqpoint{0.000000in}{0.000000in}}{\pgfqpoint{0.000000in}{0.027778in}}{%
\pgfpathmoveto{\pgfqpoint{0.000000in}{0.000000in}}%
\pgfpathlineto{\pgfqpoint{0.000000in}{0.027778in}}%
\pgfusepath{stroke,fill}%
}%
\begin{pgfscope}%
\pgfsys@transformshift{2.041485in}{0.352669in}%
\pgfsys@useobject{currentmarker}{}%
\end{pgfscope}%
\end{pgfscope}%
\begin{pgfscope}%
\pgfsetbuttcap%
\pgfsetroundjoin%
\definecolor{currentfill}{rgb}{0.000000,0.000000,0.000000}%
\pgfsetfillcolor{currentfill}%
\pgfsetlinewidth{0.301125pt}%
\definecolor{currentstroke}{rgb}{0.000000,0.000000,0.000000}%
\pgfsetstrokecolor{currentstroke}%
\pgfsetdash{}{0pt}%
\pgfsys@defobject{currentmarker}{\pgfqpoint{0.000000in}{0.000000in}}{\pgfqpoint{0.000000in}{0.027778in}}{%
\pgfpathmoveto{\pgfqpoint{0.000000in}{0.000000in}}%
\pgfpathlineto{\pgfqpoint{0.000000in}{0.027778in}}%
\pgfusepath{stroke,fill}%
}%
\begin{pgfscope}%
\pgfsys@transformshift{2.141929in}{0.352669in}%
\pgfsys@useobject{currentmarker}{}%
\end{pgfscope}%
\end{pgfscope}%
\begin{pgfscope}%
\pgfsetbuttcap%
\pgfsetroundjoin%
\definecolor{currentfill}{rgb}{0.000000,0.000000,0.000000}%
\pgfsetfillcolor{currentfill}%
\pgfsetlinewidth{0.301125pt}%
\definecolor{currentstroke}{rgb}{0.000000,0.000000,0.000000}%
\pgfsetstrokecolor{currentstroke}%
\pgfsetdash{}{0pt}%
\pgfsys@defobject{currentmarker}{\pgfqpoint{0.000000in}{0.000000in}}{\pgfqpoint{0.000000in}{0.027778in}}{%
\pgfpathmoveto{\pgfqpoint{0.000000in}{0.000000in}}%
\pgfpathlineto{\pgfqpoint{0.000000in}{0.027778in}}%
\pgfusepath{stroke,fill}%
}%
\begin{pgfscope}%
\pgfsys@transformshift{2.342818in}{0.352669in}%
\pgfsys@useobject{currentmarker}{}%
\end{pgfscope}%
\end{pgfscope}%
\begin{pgfscope}%
\pgfsetbuttcap%
\pgfsetroundjoin%
\definecolor{currentfill}{rgb}{0.000000,0.000000,0.000000}%
\pgfsetfillcolor{currentfill}%
\pgfsetlinewidth{0.301125pt}%
\definecolor{currentstroke}{rgb}{0.000000,0.000000,0.000000}%
\pgfsetstrokecolor{currentstroke}%
\pgfsetdash{}{0pt}%
\pgfsys@defobject{currentmarker}{\pgfqpoint{0.000000in}{0.000000in}}{\pgfqpoint{0.000000in}{0.027778in}}{%
\pgfpathmoveto{\pgfqpoint{0.000000in}{0.000000in}}%
\pgfpathlineto{\pgfqpoint{0.000000in}{0.027778in}}%
\pgfusepath{stroke,fill}%
}%
\begin{pgfscope}%
\pgfsys@transformshift{2.443262in}{0.352669in}%
\pgfsys@useobject{currentmarker}{}%
\end{pgfscope}%
\end{pgfscope}%
\begin{pgfscope}%
\pgfsetbuttcap%
\pgfsetroundjoin%
\definecolor{currentfill}{rgb}{0.000000,0.000000,0.000000}%
\pgfsetfillcolor{currentfill}%
\pgfsetlinewidth{0.301125pt}%
\definecolor{currentstroke}{rgb}{0.000000,0.000000,0.000000}%
\pgfsetstrokecolor{currentstroke}%
\pgfsetdash{}{0pt}%
\pgfsys@defobject{currentmarker}{\pgfqpoint{0.000000in}{0.000000in}}{\pgfqpoint{0.000000in}{0.027778in}}{%
\pgfpathmoveto{\pgfqpoint{0.000000in}{0.000000in}}%
\pgfpathlineto{\pgfqpoint{0.000000in}{0.027778in}}%
\pgfusepath{stroke,fill}%
}%
\begin{pgfscope}%
\pgfsys@transformshift{2.543707in}{0.352669in}%
\pgfsys@useobject{currentmarker}{}%
\end{pgfscope}%
\end{pgfscope}%
\begin{pgfscope}%
\pgfsetbuttcap%
\pgfsetroundjoin%
\definecolor{currentfill}{rgb}{0.000000,0.000000,0.000000}%
\pgfsetfillcolor{currentfill}%
\pgfsetlinewidth{0.301125pt}%
\definecolor{currentstroke}{rgb}{0.000000,0.000000,0.000000}%
\pgfsetstrokecolor{currentstroke}%
\pgfsetdash{}{0pt}%
\pgfsys@defobject{currentmarker}{\pgfqpoint{0.000000in}{0.000000in}}{\pgfqpoint{0.000000in}{0.027778in}}{%
\pgfpathmoveto{\pgfqpoint{0.000000in}{0.000000in}}%
\pgfpathlineto{\pgfqpoint{0.000000in}{0.027778in}}%
\pgfusepath{stroke,fill}%
}%
\begin{pgfscope}%
\pgfsys@transformshift{2.744595in}{0.352669in}%
\pgfsys@useobject{currentmarker}{}%
\end{pgfscope}%
\end{pgfscope}%
\begin{pgfscope}%
\pgfsetbuttcap%
\pgfsetroundjoin%
\definecolor{currentfill}{rgb}{0.000000,0.000000,0.000000}%
\pgfsetfillcolor{currentfill}%
\pgfsetlinewidth{0.301125pt}%
\definecolor{currentstroke}{rgb}{0.000000,0.000000,0.000000}%
\pgfsetstrokecolor{currentstroke}%
\pgfsetdash{}{0pt}%
\pgfsys@defobject{currentmarker}{\pgfqpoint{0.000000in}{0.000000in}}{\pgfqpoint{0.000000in}{0.027778in}}{%
\pgfpathmoveto{\pgfqpoint{0.000000in}{0.000000in}}%
\pgfpathlineto{\pgfqpoint{0.000000in}{0.027778in}}%
\pgfusepath{stroke,fill}%
}%
\begin{pgfscope}%
\pgfsys@transformshift{2.845039in}{0.352669in}%
\pgfsys@useobject{currentmarker}{}%
\end{pgfscope}%
\end{pgfscope}%
\begin{pgfscope}%
\pgfsetbuttcap%
\pgfsetroundjoin%
\definecolor{currentfill}{rgb}{0.000000,0.000000,0.000000}%
\pgfsetfillcolor{currentfill}%
\pgfsetlinewidth{0.301125pt}%
\definecolor{currentstroke}{rgb}{0.000000,0.000000,0.000000}%
\pgfsetstrokecolor{currentstroke}%
\pgfsetdash{}{0pt}%
\pgfsys@defobject{currentmarker}{\pgfqpoint{0.000000in}{0.000000in}}{\pgfqpoint{0.000000in}{0.027778in}}{%
\pgfpathmoveto{\pgfqpoint{0.000000in}{0.000000in}}%
\pgfpathlineto{\pgfqpoint{0.000000in}{0.027778in}}%
\pgfusepath{stroke,fill}%
}%
\begin{pgfscope}%
\pgfsys@transformshift{2.945484in}{0.352669in}%
\pgfsys@useobject{currentmarker}{}%
\end{pgfscope}%
\end{pgfscope}%
\begin{pgfscope}%
\pgfsetbuttcap%
\pgfsetroundjoin%
\definecolor{currentfill}{rgb}{0.000000,0.000000,0.000000}%
\pgfsetfillcolor{currentfill}%
\pgfsetlinewidth{0.301125pt}%
\definecolor{currentstroke}{rgb}{0.000000,0.000000,0.000000}%
\pgfsetstrokecolor{currentstroke}%
\pgfsetdash{}{0pt}%
\pgfsys@defobject{currentmarker}{\pgfqpoint{0.000000in}{0.000000in}}{\pgfqpoint{0.000000in}{0.027778in}}{%
\pgfpathmoveto{\pgfqpoint{0.000000in}{0.000000in}}%
\pgfpathlineto{\pgfqpoint{0.000000in}{0.027778in}}%
\pgfusepath{stroke,fill}%
}%
\begin{pgfscope}%
\pgfsys@transformshift{3.146372in}{0.352669in}%
\pgfsys@useobject{currentmarker}{}%
\end{pgfscope}%
\end{pgfscope}%
\begin{pgfscope}%
\pgfsetbuttcap%
\pgfsetroundjoin%
\definecolor{currentfill}{rgb}{0.000000,0.000000,0.000000}%
\pgfsetfillcolor{currentfill}%
\pgfsetlinewidth{0.301125pt}%
\definecolor{currentstroke}{rgb}{0.000000,0.000000,0.000000}%
\pgfsetstrokecolor{currentstroke}%
\pgfsetdash{}{0pt}%
\pgfsys@defobject{currentmarker}{\pgfqpoint{0.000000in}{0.000000in}}{\pgfqpoint{0.000000in}{0.027778in}}{%
\pgfpathmoveto{\pgfqpoint{0.000000in}{0.000000in}}%
\pgfpathlineto{\pgfqpoint{0.000000in}{0.027778in}}%
\pgfusepath{stroke,fill}%
}%
\begin{pgfscope}%
\pgfsys@transformshift{3.246817in}{0.352669in}%
\pgfsys@useobject{currentmarker}{}%
\end{pgfscope}%
\end{pgfscope}%
\begin{pgfscope}%
\pgfsetbuttcap%
\pgfsetroundjoin%
\definecolor{currentfill}{rgb}{0.000000,0.000000,0.000000}%
\pgfsetfillcolor{currentfill}%
\pgfsetlinewidth{0.301125pt}%
\definecolor{currentstroke}{rgb}{0.000000,0.000000,0.000000}%
\pgfsetstrokecolor{currentstroke}%
\pgfsetdash{}{0pt}%
\pgfsys@defobject{currentmarker}{\pgfqpoint{0.000000in}{0.000000in}}{\pgfqpoint{0.000000in}{0.027778in}}{%
\pgfpathmoveto{\pgfqpoint{0.000000in}{0.000000in}}%
\pgfpathlineto{\pgfqpoint{0.000000in}{0.027778in}}%
\pgfusepath{stroke,fill}%
}%
\begin{pgfscope}%
\pgfsys@transformshift{3.347261in}{0.352669in}%
\pgfsys@useobject{currentmarker}{}%
\end{pgfscope}%
\end{pgfscope}%
\begin{pgfscope}%
\pgfsetbuttcap%
\pgfsetroundjoin%
\definecolor{currentfill}{rgb}{0.000000,0.000000,0.000000}%
\pgfsetfillcolor{currentfill}%
\pgfsetlinewidth{0.301125pt}%
\definecolor{currentstroke}{rgb}{0.000000,0.000000,0.000000}%
\pgfsetstrokecolor{currentstroke}%
\pgfsetdash{}{0pt}%
\pgfsys@defobject{currentmarker}{\pgfqpoint{0.000000in}{0.000000in}}{\pgfqpoint{0.000000in}{0.027778in}}{%
\pgfpathmoveto{\pgfqpoint{0.000000in}{0.000000in}}%
\pgfpathlineto{\pgfqpoint{0.000000in}{0.027778in}}%
\pgfusepath{stroke,fill}%
}%
\begin{pgfscope}%
\pgfsys@transformshift{3.548149in}{0.352669in}%
\pgfsys@useobject{currentmarker}{}%
\end{pgfscope}%
\end{pgfscope}%
\begin{pgfscope}%
\pgfsetbuttcap%
\pgfsetroundjoin%
\definecolor{currentfill}{rgb}{0.000000,0.000000,0.000000}%
\pgfsetfillcolor{currentfill}%
\pgfsetlinewidth{0.301125pt}%
\definecolor{currentstroke}{rgb}{0.000000,0.000000,0.000000}%
\pgfsetstrokecolor{currentstroke}%
\pgfsetdash{}{0pt}%
\pgfsys@defobject{currentmarker}{\pgfqpoint{0.000000in}{0.000000in}}{\pgfqpoint{0.000000in}{0.027778in}}{%
\pgfpathmoveto{\pgfqpoint{0.000000in}{0.000000in}}%
\pgfpathlineto{\pgfqpoint{0.000000in}{0.027778in}}%
\pgfusepath{stroke,fill}%
}%
\begin{pgfscope}%
\pgfsys@transformshift{3.648594in}{0.352669in}%
\pgfsys@useobject{currentmarker}{}%
\end{pgfscope}%
\end{pgfscope}%
\begin{pgfscope}%
\pgfsetbuttcap%
\pgfsetroundjoin%
\definecolor{currentfill}{rgb}{0.000000,0.000000,0.000000}%
\pgfsetfillcolor{currentfill}%
\pgfsetlinewidth{0.301125pt}%
\definecolor{currentstroke}{rgb}{0.000000,0.000000,0.000000}%
\pgfsetstrokecolor{currentstroke}%
\pgfsetdash{}{0pt}%
\pgfsys@defobject{currentmarker}{\pgfqpoint{0.000000in}{0.000000in}}{\pgfqpoint{0.000000in}{0.027778in}}{%
\pgfpathmoveto{\pgfqpoint{0.000000in}{0.000000in}}%
\pgfpathlineto{\pgfqpoint{0.000000in}{0.027778in}}%
\pgfusepath{stroke,fill}%
}%
\begin{pgfscope}%
\pgfsys@transformshift{3.749038in}{0.352669in}%
\pgfsys@useobject{currentmarker}{}%
\end{pgfscope}%
\end{pgfscope}%
\begin{pgfscope}%
\definecolor{textcolor}{rgb}{0.000000,0.000000,0.000000}%
\pgfsetstrokecolor{textcolor}%
\pgfsetfillcolor{textcolor}%
\pgftext[x=2.149000in,y=0.140972in,,top]{\color{textcolor}{\sffamily\fontsize{9.000000}{10.800000}\selectfont\catcode`\^=\active\def^{\ifmmode\sp\else\^{}\fi}\catcode`\%=\active\def%{\%}$\lambda$}}%
\end{pgfscope}%
\begin{pgfscope}%
\pgfsetbuttcap%
\pgfsetroundjoin%
\definecolor{currentfill}{rgb}{0.000000,0.000000,0.000000}%
\pgfsetfillcolor{currentfill}%
\pgfsetlinewidth{0.501875pt}%
\definecolor{currentstroke}{rgb}{0.000000,0.000000,0.000000}%
\pgfsetstrokecolor{currentstroke}%
\pgfsetdash{}{0pt}%
\pgfsys@defobject{currentmarker}{\pgfqpoint{0.000000in}{0.000000in}}{\pgfqpoint{0.041667in}{0.000000in}}{%
\pgfpathmoveto{\pgfqpoint{0.000000in}{0.000000in}}%
\pgfpathlineto{\pgfqpoint{0.041667in}{0.000000in}}%
\pgfusepath{stroke,fill}%
}%
\begin{pgfscope}%
\pgfsys@transformshift{0.484000in}{0.431957in}%
\pgfsys@useobject{currentmarker}{}%
\end{pgfscope}%
\end{pgfscope}%
\begin{pgfscope}%
\definecolor{textcolor}{rgb}{0.000000,0.000000,0.000000}%
\pgfsetstrokecolor{textcolor}%
\pgfsetfillcolor{textcolor}%
\pgftext[x=0.196694in, y=0.389748in, left, base]{\color{textcolor}{\sffamily\fontsize{8.000000}{9.600000}\selectfont\catcode`\^=\active\def^{\ifmmode\sp\else\^{}\fi}\catcode`\%=\active\def%{\%}$\mathdefault{10^{-2}}$}}%
\end{pgfscope}%
\begin{pgfscope}%
\pgfsetbuttcap%
\pgfsetroundjoin%
\definecolor{currentfill}{rgb}{0.000000,0.000000,0.000000}%
\pgfsetfillcolor{currentfill}%
\pgfsetlinewidth{0.501875pt}%
\definecolor{currentstroke}{rgb}{0.000000,0.000000,0.000000}%
\pgfsetstrokecolor{currentstroke}%
\pgfsetdash{}{0pt}%
\pgfsys@defobject{currentmarker}{\pgfqpoint{0.000000in}{0.000000in}}{\pgfqpoint{0.041667in}{0.000000in}}{%
\pgfpathmoveto{\pgfqpoint{0.000000in}{0.000000in}}%
\pgfpathlineto{\pgfqpoint{0.041667in}{0.000000in}}%
\pgfusepath{stroke,fill}%
}%
\begin{pgfscope}%
\pgfsys@transformshift{0.484000in}{0.955972in}%
\pgfsys@useobject{currentmarker}{}%
\end{pgfscope}%
\end{pgfscope}%
\begin{pgfscope}%
\definecolor{textcolor}{rgb}{0.000000,0.000000,0.000000}%
\pgfsetstrokecolor{textcolor}%
\pgfsetfillcolor{textcolor}%
\pgftext[x=0.196527in, y=0.913763in, left, base]{\color{textcolor}{\sffamily\fontsize{8.000000}{9.600000}\selectfont\catcode`\^=\active\def^{\ifmmode\sp\else\^{}\fi}\catcode`\%=\active\def%{\%}$\mathdefault{10^{-1}}$}}%
\end{pgfscope}%
\begin{pgfscope}%
\pgfsetbuttcap%
\pgfsetroundjoin%
\definecolor{currentfill}{rgb}{0.000000,0.000000,0.000000}%
\pgfsetfillcolor{currentfill}%
\pgfsetlinewidth{0.501875pt}%
\definecolor{currentstroke}{rgb}{0.000000,0.000000,0.000000}%
\pgfsetstrokecolor{currentstroke}%
\pgfsetdash{}{0pt}%
\pgfsys@defobject{currentmarker}{\pgfqpoint{0.000000in}{0.000000in}}{\pgfqpoint{0.041667in}{0.000000in}}{%
\pgfpathmoveto{\pgfqpoint{0.000000in}{0.000000in}}%
\pgfpathlineto{\pgfqpoint{0.041667in}{0.000000in}}%
\pgfusepath{stroke,fill}%
}%
\begin{pgfscope}%
\pgfsys@transformshift{0.484000in}{1.479988in}%
\pgfsys@useobject{currentmarker}{}%
\end{pgfscope}%
\end{pgfscope}%
\begin{pgfscope}%
\definecolor{textcolor}{rgb}{0.000000,0.000000,0.000000}%
\pgfsetstrokecolor{textcolor}%
\pgfsetfillcolor{textcolor}%
\pgftext[x=0.259444in, y=1.437779in, left, base]{\color{textcolor}{\sffamily\fontsize{8.000000}{9.600000}\selectfont\catcode`\^=\active\def^{\ifmmode\sp\else\^{}\fi}\catcode`\%=\active\def%{\%}$\mathdefault{10^{0}}$}}%
\end{pgfscope}%
\begin{pgfscope}%
\pgfsetbuttcap%
\pgfsetroundjoin%
\definecolor{currentfill}{rgb}{0.000000,0.000000,0.000000}%
\pgfsetfillcolor{currentfill}%
\pgfsetlinewidth{0.301125pt}%
\definecolor{currentstroke}{rgb}{0.000000,0.000000,0.000000}%
\pgfsetstrokecolor{currentstroke}%
\pgfsetdash{}{0pt}%
\pgfsys@defobject{currentmarker}{\pgfqpoint{0.000000in}{0.000000in}}{\pgfqpoint{0.027778in}{0.000000in}}{%
\pgfpathmoveto{\pgfqpoint{0.000000in}{0.000000in}}%
\pgfpathlineto{\pgfqpoint{0.027778in}{0.000000in}}%
\pgfusepath{stroke,fill}%
}%
\begin{pgfscope}%
\pgfsys@transformshift{0.484000in}{0.381175in}%
\pgfsys@useobject{currentmarker}{}%
\end{pgfscope}%
\end{pgfscope}%
\begin{pgfscope}%
\pgfsetbuttcap%
\pgfsetroundjoin%
\definecolor{currentfill}{rgb}{0.000000,0.000000,0.000000}%
\pgfsetfillcolor{currentfill}%
\pgfsetlinewidth{0.301125pt}%
\definecolor{currentstroke}{rgb}{0.000000,0.000000,0.000000}%
\pgfsetstrokecolor{currentstroke}%
\pgfsetdash{}{0pt}%
\pgfsys@defobject{currentmarker}{\pgfqpoint{0.000000in}{0.000000in}}{\pgfqpoint{0.027778in}{0.000000in}}{%
\pgfpathmoveto{\pgfqpoint{0.000000in}{0.000000in}}%
\pgfpathlineto{\pgfqpoint{0.027778in}{0.000000in}}%
\pgfusepath{stroke,fill}%
}%
\begin{pgfscope}%
\pgfsys@transformshift{0.484000in}{0.407979in}%
\pgfsys@useobject{currentmarker}{}%
\end{pgfscope}%
\end{pgfscope}%
\begin{pgfscope}%
\pgfsetbuttcap%
\pgfsetroundjoin%
\definecolor{currentfill}{rgb}{0.000000,0.000000,0.000000}%
\pgfsetfillcolor{currentfill}%
\pgfsetlinewidth{0.301125pt}%
\definecolor{currentstroke}{rgb}{0.000000,0.000000,0.000000}%
\pgfsetstrokecolor{currentstroke}%
\pgfsetdash{}{0pt}%
\pgfsys@defobject{currentmarker}{\pgfqpoint{0.000000in}{0.000000in}}{\pgfqpoint{0.027778in}{0.000000in}}{%
\pgfpathmoveto{\pgfqpoint{0.000000in}{0.000000in}}%
\pgfpathlineto{\pgfqpoint{0.027778in}{0.000000in}}%
\pgfusepath{stroke,fill}%
}%
\begin{pgfscope}%
\pgfsys@transformshift{0.484000in}{0.589701in}%
\pgfsys@useobject{currentmarker}{}%
\end{pgfscope}%
\end{pgfscope}%
\begin{pgfscope}%
\pgfsetbuttcap%
\pgfsetroundjoin%
\definecolor{currentfill}{rgb}{0.000000,0.000000,0.000000}%
\pgfsetfillcolor{currentfill}%
\pgfsetlinewidth{0.301125pt}%
\definecolor{currentstroke}{rgb}{0.000000,0.000000,0.000000}%
\pgfsetstrokecolor{currentstroke}%
\pgfsetdash{}{0pt}%
\pgfsys@defobject{currentmarker}{\pgfqpoint{0.000000in}{0.000000in}}{\pgfqpoint{0.027778in}{0.000000in}}{%
\pgfpathmoveto{\pgfqpoint{0.000000in}{0.000000in}}%
\pgfpathlineto{\pgfqpoint{0.027778in}{0.000000in}}%
\pgfusepath{stroke,fill}%
}%
\begin{pgfscope}%
\pgfsys@transformshift{0.484000in}{0.681976in}%
\pgfsys@useobject{currentmarker}{}%
\end{pgfscope}%
\end{pgfscope}%
\begin{pgfscope}%
\pgfsetbuttcap%
\pgfsetroundjoin%
\definecolor{currentfill}{rgb}{0.000000,0.000000,0.000000}%
\pgfsetfillcolor{currentfill}%
\pgfsetlinewidth{0.301125pt}%
\definecolor{currentstroke}{rgb}{0.000000,0.000000,0.000000}%
\pgfsetstrokecolor{currentstroke}%
\pgfsetdash{}{0pt}%
\pgfsys@defobject{currentmarker}{\pgfqpoint{0.000000in}{0.000000in}}{\pgfqpoint{0.027778in}{0.000000in}}{%
\pgfpathmoveto{\pgfqpoint{0.000000in}{0.000000in}}%
\pgfpathlineto{\pgfqpoint{0.027778in}{0.000000in}}%
\pgfusepath{stroke,fill}%
}%
\begin{pgfscope}%
\pgfsys@transformshift{0.484000in}{0.747446in}%
\pgfsys@useobject{currentmarker}{}%
\end{pgfscope}%
\end{pgfscope}%
\begin{pgfscope}%
\pgfsetbuttcap%
\pgfsetroundjoin%
\definecolor{currentfill}{rgb}{0.000000,0.000000,0.000000}%
\pgfsetfillcolor{currentfill}%
\pgfsetlinewidth{0.301125pt}%
\definecolor{currentstroke}{rgb}{0.000000,0.000000,0.000000}%
\pgfsetstrokecolor{currentstroke}%
\pgfsetdash{}{0pt}%
\pgfsys@defobject{currentmarker}{\pgfqpoint{0.000000in}{0.000000in}}{\pgfqpoint{0.027778in}{0.000000in}}{%
\pgfpathmoveto{\pgfqpoint{0.000000in}{0.000000in}}%
\pgfpathlineto{\pgfqpoint{0.027778in}{0.000000in}}%
\pgfusepath{stroke,fill}%
}%
\begin{pgfscope}%
\pgfsys@transformshift{0.484000in}{0.798228in}%
\pgfsys@useobject{currentmarker}{}%
\end{pgfscope}%
\end{pgfscope}%
\begin{pgfscope}%
\pgfsetbuttcap%
\pgfsetroundjoin%
\definecolor{currentfill}{rgb}{0.000000,0.000000,0.000000}%
\pgfsetfillcolor{currentfill}%
\pgfsetlinewidth{0.301125pt}%
\definecolor{currentstroke}{rgb}{0.000000,0.000000,0.000000}%
\pgfsetstrokecolor{currentstroke}%
\pgfsetdash{}{0pt}%
\pgfsys@defobject{currentmarker}{\pgfqpoint{0.000000in}{0.000000in}}{\pgfqpoint{0.027778in}{0.000000in}}{%
\pgfpathmoveto{\pgfqpoint{0.000000in}{0.000000in}}%
\pgfpathlineto{\pgfqpoint{0.027778in}{0.000000in}}%
\pgfusepath{stroke,fill}%
}%
\begin{pgfscope}%
\pgfsys@transformshift{0.484000in}{0.839720in}%
\pgfsys@useobject{currentmarker}{}%
\end{pgfscope}%
\end{pgfscope}%
\begin{pgfscope}%
\pgfsetbuttcap%
\pgfsetroundjoin%
\definecolor{currentfill}{rgb}{0.000000,0.000000,0.000000}%
\pgfsetfillcolor{currentfill}%
\pgfsetlinewidth{0.301125pt}%
\definecolor{currentstroke}{rgb}{0.000000,0.000000,0.000000}%
\pgfsetstrokecolor{currentstroke}%
\pgfsetdash{}{0pt}%
\pgfsys@defobject{currentmarker}{\pgfqpoint{0.000000in}{0.000000in}}{\pgfqpoint{0.027778in}{0.000000in}}{%
\pgfpathmoveto{\pgfqpoint{0.000000in}{0.000000in}}%
\pgfpathlineto{\pgfqpoint{0.027778in}{0.000000in}}%
\pgfusepath{stroke,fill}%
}%
\begin{pgfscope}%
\pgfsys@transformshift{0.484000in}{0.874801in}%
\pgfsys@useobject{currentmarker}{}%
\end{pgfscope}%
\end{pgfscope}%
\begin{pgfscope}%
\pgfsetbuttcap%
\pgfsetroundjoin%
\definecolor{currentfill}{rgb}{0.000000,0.000000,0.000000}%
\pgfsetfillcolor{currentfill}%
\pgfsetlinewidth{0.301125pt}%
\definecolor{currentstroke}{rgb}{0.000000,0.000000,0.000000}%
\pgfsetstrokecolor{currentstroke}%
\pgfsetdash{}{0pt}%
\pgfsys@defobject{currentmarker}{\pgfqpoint{0.000000in}{0.000000in}}{\pgfqpoint{0.027778in}{0.000000in}}{%
\pgfpathmoveto{\pgfqpoint{0.000000in}{0.000000in}}%
\pgfpathlineto{\pgfqpoint{0.027778in}{0.000000in}}%
\pgfusepath{stroke,fill}%
}%
\begin{pgfscope}%
\pgfsys@transformshift{0.484000in}{0.905190in}%
\pgfsys@useobject{currentmarker}{}%
\end{pgfscope}%
\end{pgfscope}%
\begin{pgfscope}%
\pgfsetbuttcap%
\pgfsetroundjoin%
\definecolor{currentfill}{rgb}{0.000000,0.000000,0.000000}%
\pgfsetfillcolor{currentfill}%
\pgfsetlinewidth{0.301125pt}%
\definecolor{currentstroke}{rgb}{0.000000,0.000000,0.000000}%
\pgfsetstrokecolor{currentstroke}%
\pgfsetdash{}{0pt}%
\pgfsys@defobject{currentmarker}{\pgfqpoint{0.000000in}{0.000000in}}{\pgfqpoint{0.027778in}{0.000000in}}{%
\pgfpathmoveto{\pgfqpoint{0.000000in}{0.000000in}}%
\pgfpathlineto{\pgfqpoint{0.027778in}{0.000000in}}%
\pgfusepath{stroke,fill}%
}%
\begin{pgfscope}%
\pgfsys@transformshift{0.484000in}{0.931995in}%
\pgfsys@useobject{currentmarker}{}%
\end{pgfscope}%
\end{pgfscope}%
\begin{pgfscope}%
\pgfsetbuttcap%
\pgfsetroundjoin%
\definecolor{currentfill}{rgb}{0.000000,0.000000,0.000000}%
\pgfsetfillcolor{currentfill}%
\pgfsetlinewidth{0.301125pt}%
\definecolor{currentstroke}{rgb}{0.000000,0.000000,0.000000}%
\pgfsetstrokecolor{currentstroke}%
\pgfsetdash{}{0pt}%
\pgfsys@defobject{currentmarker}{\pgfqpoint{0.000000in}{0.000000in}}{\pgfqpoint{0.027778in}{0.000000in}}{%
\pgfpathmoveto{\pgfqpoint{0.000000in}{0.000000in}}%
\pgfpathlineto{\pgfqpoint{0.027778in}{0.000000in}}%
\pgfusepath{stroke,fill}%
}%
\begin{pgfscope}%
\pgfsys@transformshift{0.484000in}{1.113717in}%
\pgfsys@useobject{currentmarker}{}%
\end{pgfscope}%
\end{pgfscope}%
\begin{pgfscope}%
\pgfsetbuttcap%
\pgfsetroundjoin%
\definecolor{currentfill}{rgb}{0.000000,0.000000,0.000000}%
\pgfsetfillcolor{currentfill}%
\pgfsetlinewidth{0.301125pt}%
\definecolor{currentstroke}{rgb}{0.000000,0.000000,0.000000}%
\pgfsetstrokecolor{currentstroke}%
\pgfsetdash{}{0pt}%
\pgfsys@defobject{currentmarker}{\pgfqpoint{0.000000in}{0.000000in}}{\pgfqpoint{0.027778in}{0.000000in}}{%
\pgfpathmoveto{\pgfqpoint{0.000000in}{0.000000in}}%
\pgfpathlineto{\pgfqpoint{0.027778in}{0.000000in}}%
\pgfusepath{stroke,fill}%
}%
\begin{pgfscope}%
\pgfsys@transformshift{0.484000in}{1.205991in}%
\pgfsys@useobject{currentmarker}{}%
\end{pgfscope}%
\end{pgfscope}%
\begin{pgfscope}%
\pgfsetbuttcap%
\pgfsetroundjoin%
\definecolor{currentfill}{rgb}{0.000000,0.000000,0.000000}%
\pgfsetfillcolor{currentfill}%
\pgfsetlinewidth{0.301125pt}%
\definecolor{currentstroke}{rgb}{0.000000,0.000000,0.000000}%
\pgfsetstrokecolor{currentstroke}%
\pgfsetdash{}{0pt}%
\pgfsys@defobject{currentmarker}{\pgfqpoint{0.000000in}{0.000000in}}{\pgfqpoint{0.027778in}{0.000000in}}{%
\pgfpathmoveto{\pgfqpoint{0.000000in}{0.000000in}}%
\pgfpathlineto{\pgfqpoint{0.027778in}{0.000000in}}%
\pgfusepath{stroke,fill}%
}%
\begin{pgfscope}%
\pgfsys@transformshift{0.484000in}{1.271461in}%
\pgfsys@useobject{currentmarker}{}%
\end{pgfscope}%
\end{pgfscope}%
\begin{pgfscope}%
\pgfsetbuttcap%
\pgfsetroundjoin%
\definecolor{currentfill}{rgb}{0.000000,0.000000,0.000000}%
\pgfsetfillcolor{currentfill}%
\pgfsetlinewidth{0.301125pt}%
\definecolor{currentstroke}{rgb}{0.000000,0.000000,0.000000}%
\pgfsetstrokecolor{currentstroke}%
\pgfsetdash{}{0pt}%
\pgfsys@defobject{currentmarker}{\pgfqpoint{0.000000in}{0.000000in}}{\pgfqpoint{0.027778in}{0.000000in}}{%
\pgfpathmoveto{\pgfqpoint{0.000000in}{0.000000in}}%
\pgfpathlineto{\pgfqpoint{0.027778in}{0.000000in}}%
\pgfusepath{stroke,fill}%
}%
\begin{pgfscope}%
\pgfsys@transformshift{0.484000in}{1.322244in}%
\pgfsys@useobject{currentmarker}{}%
\end{pgfscope}%
\end{pgfscope}%
\begin{pgfscope}%
\pgfsetbuttcap%
\pgfsetroundjoin%
\definecolor{currentfill}{rgb}{0.000000,0.000000,0.000000}%
\pgfsetfillcolor{currentfill}%
\pgfsetlinewidth{0.301125pt}%
\definecolor{currentstroke}{rgb}{0.000000,0.000000,0.000000}%
\pgfsetstrokecolor{currentstroke}%
\pgfsetdash{}{0pt}%
\pgfsys@defobject{currentmarker}{\pgfqpoint{0.000000in}{0.000000in}}{\pgfqpoint{0.027778in}{0.000000in}}{%
\pgfpathmoveto{\pgfqpoint{0.000000in}{0.000000in}}%
\pgfpathlineto{\pgfqpoint{0.027778in}{0.000000in}}%
\pgfusepath{stroke,fill}%
}%
\begin{pgfscope}%
\pgfsys@transformshift{0.484000in}{1.363736in}%
\pgfsys@useobject{currentmarker}{}%
\end{pgfscope}%
\end{pgfscope}%
\begin{pgfscope}%
\pgfsetbuttcap%
\pgfsetroundjoin%
\definecolor{currentfill}{rgb}{0.000000,0.000000,0.000000}%
\pgfsetfillcolor{currentfill}%
\pgfsetlinewidth{0.301125pt}%
\definecolor{currentstroke}{rgb}{0.000000,0.000000,0.000000}%
\pgfsetstrokecolor{currentstroke}%
\pgfsetdash{}{0pt}%
\pgfsys@defobject{currentmarker}{\pgfqpoint{0.000000in}{0.000000in}}{\pgfqpoint{0.027778in}{0.000000in}}{%
\pgfpathmoveto{\pgfqpoint{0.000000in}{0.000000in}}%
\pgfpathlineto{\pgfqpoint{0.027778in}{0.000000in}}%
\pgfusepath{stroke,fill}%
}%
\begin{pgfscope}%
\pgfsys@transformshift{0.484000in}{1.398817in}%
\pgfsys@useobject{currentmarker}{}%
\end{pgfscope}%
\end{pgfscope}%
\begin{pgfscope}%
\pgfsetbuttcap%
\pgfsetroundjoin%
\definecolor{currentfill}{rgb}{0.000000,0.000000,0.000000}%
\pgfsetfillcolor{currentfill}%
\pgfsetlinewidth{0.301125pt}%
\definecolor{currentstroke}{rgb}{0.000000,0.000000,0.000000}%
\pgfsetstrokecolor{currentstroke}%
\pgfsetdash{}{0pt}%
\pgfsys@defobject{currentmarker}{\pgfqpoint{0.000000in}{0.000000in}}{\pgfqpoint{0.027778in}{0.000000in}}{%
\pgfpathmoveto{\pgfqpoint{0.000000in}{0.000000in}}%
\pgfpathlineto{\pgfqpoint{0.027778in}{0.000000in}}%
\pgfusepath{stroke,fill}%
}%
\begin{pgfscope}%
\pgfsys@transformshift{0.484000in}{1.429206in}%
\pgfsys@useobject{currentmarker}{}%
\end{pgfscope}%
\end{pgfscope}%
\begin{pgfscope}%
\pgfsetbuttcap%
\pgfsetroundjoin%
\definecolor{currentfill}{rgb}{0.000000,0.000000,0.000000}%
\pgfsetfillcolor{currentfill}%
\pgfsetlinewidth{0.301125pt}%
\definecolor{currentstroke}{rgb}{0.000000,0.000000,0.000000}%
\pgfsetstrokecolor{currentstroke}%
\pgfsetdash{}{0pt}%
\pgfsys@defobject{currentmarker}{\pgfqpoint{0.000000in}{0.000000in}}{\pgfqpoint{0.027778in}{0.000000in}}{%
\pgfpathmoveto{\pgfqpoint{0.000000in}{0.000000in}}%
\pgfpathlineto{\pgfqpoint{0.027778in}{0.000000in}}%
\pgfusepath{stroke,fill}%
}%
\begin{pgfscope}%
\pgfsys@transformshift{0.484000in}{1.456010in}%
\pgfsys@useobject{currentmarker}{}%
\end{pgfscope}%
\end{pgfscope}%
\begin{pgfscope}%
\pgfsetbuttcap%
\pgfsetroundjoin%
\definecolor{currentfill}{rgb}{0.000000,0.000000,0.000000}%
\pgfsetfillcolor{currentfill}%
\pgfsetlinewidth{0.301125pt}%
\definecolor{currentstroke}{rgb}{0.000000,0.000000,0.000000}%
\pgfsetstrokecolor{currentstroke}%
\pgfsetdash{}{0pt}%
\pgfsys@defobject{currentmarker}{\pgfqpoint{0.000000in}{0.000000in}}{\pgfqpoint{0.027778in}{0.000000in}}{%
\pgfpathmoveto{\pgfqpoint{0.000000in}{0.000000in}}%
\pgfpathlineto{\pgfqpoint{0.027778in}{0.000000in}}%
\pgfusepath{stroke,fill}%
}%
\begin{pgfscope}%
\pgfsys@transformshift{0.484000in}{1.637732in}%
\pgfsys@useobject{currentmarker}{}%
\end{pgfscope}%
\end{pgfscope}%
\begin{pgfscope}%
\pgfsetbuttcap%
\pgfsetroundjoin%
\definecolor{currentfill}{rgb}{0.000000,0.000000,0.000000}%
\pgfsetfillcolor{currentfill}%
\pgfsetlinewidth{0.301125pt}%
\definecolor{currentstroke}{rgb}{0.000000,0.000000,0.000000}%
\pgfsetstrokecolor{currentstroke}%
\pgfsetdash{}{0pt}%
\pgfsys@defobject{currentmarker}{\pgfqpoint{0.000000in}{0.000000in}}{\pgfqpoint{0.027778in}{0.000000in}}{%
\pgfpathmoveto{\pgfqpoint{0.000000in}{0.000000in}}%
\pgfpathlineto{\pgfqpoint{0.027778in}{0.000000in}}%
\pgfusepath{stroke,fill}%
}%
\begin{pgfscope}%
\pgfsys@transformshift{0.484000in}{1.730007in}%
\pgfsys@useobject{currentmarker}{}%
\end{pgfscope}%
\end{pgfscope}%
\begin{pgfscope}%
\pgfsetbuttcap%
\pgfsetroundjoin%
\definecolor{currentfill}{rgb}{0.000000,0.000000,0.000000}%
\pgfsetfillcolor{currentfill}%
\pgfsetlinewidth{0.301125pt}%
\definecolor{currentstroke}{rgb}{0.000000,0.000000,0.000000}%
\pgfsetstrokecolor{currentstroke}%
\pgfsetdash{}{0pt}%
\pgfsys@defobject{currentmarker}{\pgfqpoint{0.000000in}{0.000000in}}{\pgfqpoint{0.027778in}{0.000000in}}{%
\pgfpathmoveto{\pgfqpoint{0.000000in}{0.000000in}}%
\pgfpathlineto{\pgfqpoint{0.027778in}{0.000000in}}%
\pgfusepath{stroke,fill}%
}%
\begin{pgfscope}%
\pgfsys@transformshift{0.484000in}{1.795477in}%
\pgfsys@useobject{currentmarker}{}%
\end{pgfscope}%
\end{pgfscope}%
\begin{pgfscope}%
\pgfsetbuttcap%
\pgfsetroundjoin%
\definecolor{currentfill}{rgb}{0.000000,0.000000,0.000000}%
\pgfsetfillcolor{currentfill}%
\pgfsetlinewidth{0.301125pt}%
\definecolor{currentstroke}{rgb}{0.000000,0.000000,0.000000}%
\pgfsetstrokecolor{currentstroke}%
\pgfsetdash{}{0pt}%
\pgfsys@defobject{currentmarker}{\pgfqpoint{0.000000in}{0.000000in}}{\pgfqpoint{0.027778in}{0.000000in}}{%
\pgfpathmoveto{\pgfqpoint{0.000000in}{0.000000in}}%
\pgfpathlineto{\pgfqpoint{0.027778in}{0.000000in}}%
\pgfusepath{stroke,fill}%
}%
\begin{pgfscope}%
\pgfsys@transformshift{0.484000in}{1.846259in}%
\pgfsys@useobject{currentmarker}{}%
\end{pgfscope}%
\end{pgfscope}%
\begin{pgfscope}%
\pgfsetbuttcap%
\pgfsetroundjoin%
\definecolor{currentfill}{rgb}{0.000000,0.000000,0.000000}%
\pgfsetfillcolor{currentfill}%
\pgfsetlinewidth{0.301125pt}%
\definecolor{currentstroke}{rgb}{0.000000,0.000000,0.000000}%
\pgfsetstrokecolor{currentstroke}%
\pgfsetdash{}{0pt}%
\pgfsys@defobject{currentmarker}{\pgfqpoint{0.000000in}{0.000000in}}{\pgfqpoint{0.027778in}{0.000000in}}{%
\pgfpathmoveto{\pgfqpoint{0.000000in}{0.000000in}}%
\pgfpathlineto{\pgfqpoint{0.027778in}{0.000000in}}%
\pgfusepath{stroke,fill}%
}%
\begin{pgfscope}%
\pgfsys@transformshift{0.484000in}{1.887751in}%
\pgfsys@useobject{currentmarker}{}%
\end{pgfscope}%
\end{pgfscope}%
\begin{pgfscope}%
\pgfsetbuttcap%
\pgfsetroundjoin%
\definecolor{currentfill}{rgb}{0.000000,0.000000,0.000000}%
\pgfsetfillcolor{currentfill}%
\pgfsetlinewidth{0.301125pt}%
\definecolor{currentstroke}{rgb}{0.000000,0.000000,0.000000}%
\pgfsetstrokecolor{currentstroke}%
\pgfsetdash{}{0pt}%
\pgfsys@defobject{currentmarker}{\pgfqpoint{0.000000in}{0.000000in}}{\pgfqpoint{0.027778in}{0.000000in}}{%
\pgfpathmoveto{\pgfqpoint{0.000000in}{0.000000in}}%
\pgfpathlineto{\pgfqpoint{0.027778in}{0.000000in}}%
\pgfusepath{stroke,fill}%
}%
\begin{pgfscope}%
\pgfsys@transformshift{0.484000in}{1.922832in}%
\pgfsys@useobject{currentmarker}{}%
\end{pgfscope}%
\end{pgfscope}%
\begin{pgfscope}%
\pgfsetbuttcap%
\pgfsetroundjoin%
\definecolor{currentfill}{rgb}{0.000000,0.000000,0.000000}%
\pgfsetfillcolor{currentfill}%
\pgfsetlinewidth{0.301125pt}%
\definecolor{currentstroke}{rgb}{0.000000,0.000000,0.000000}%
\pgfsetstrokecolor{currentstroke}%
\pgfsetdash{}{0pt}%
\pgfsys@defobject{currentmarker}{\pgfqpoint{0.000000in}{0.000000in}}{\pgfqpoint{0.027778in}{0.000000in}}{%
\pgfpathmoveto{\pgfqpoint{0.000000in}{0.000000in}}%
\pgfpathlineto{\pgfqpoint{0.027778in}{0.000000in}}%
\pgfusepath{stroke,fill}%
}%
\begin{pgfscope}%
\pgfsys@transformshift{0.484000in}{1.953221in}%
\pgfsys@useobject{currentmarker}{}%
\end{pgfscope}%
\end{pgfscope}%
\begin{pgfscope}%
\definecolor{textcolor}{rgb}{0.000000,0.000000,0.000000}%
\pgfsetstrokecolor{textcolor}%
\pgfsetfillcolor{textcolor}%
\pgftext[x=0.140972in,y=1.162184in,,bottom,rotate=90.000000]{\color{textcolor}{\sffamily\fontsize{9.000000}{10.800000}\selectfont\catcode`\^=\active\def^{\ifmmode\sp\else\^{}\fi}\catcode`\%=\active\def%{\%}$h$}}%
\end{pgfscope}%
\begin{pgfscope}%
\pgfpathrectangle{\pgfqpoint{0.484000in}{0.352669in}}{\pgfqpoint{3.330000in}{1.619031in}}%
\pgfusepath{clip}%
\pgfsetrectcap%
\pgfsetroundjoin%
\pgfsetlinewidth{0.803000pt}%
\definecolor{currentstroke}{rgb}{0.768627,0.556863,0.211765}%
\pgfsetstrokecolor{currentstroke}%
\pgfsetdash{}{0pt}%
\pgfpathmoveto{\pgfqpoint{0.635363in}{0.426261in}}%
\pgfpathlineto{\pgfqpoint{0.636441in}{0.950277in}}%
\pgfpathlineto{\pgfqpoint{0.647217in}{1.474292in}}%
\pgfpathlineto{\pgfqpoint{0.684364in}{1.707961in}}%
\pgfpathlineto{\pgfqpoint{0.737873in}{1.697810in}}%
\pgfpathlineto{\pgfqpoint{0.788552in}{1.683104in}}%
\pgfpathlineto{\pgfqpoint{0.880629in}{1.653964in}}%
\pgfpathlineto{\pgfqpoint{0.961766in}{1.625553in}}%
\pgfpathlineto{\pgfqpoint{1.033484in}{1.597802in}}%
\pgfpathlineto{\pgfqpoint{1.097044in}{1.570591in}}%
\pgfpathlineto{\pgfqpoint{1.153486in}{1.543719in}}%
\pgfpathlineto{\pgfqpoint{1.203650in}{1.516889in}}%
\pgfpathlineto{\pgfqpoint{1.226595in}{1.503376in}}%
\pgfpathlineto{\pgfqpoint{1.248211in}{1.489744in}}%
\pgfpathlineto{\pgfqpoint{1.268564in}{1.475969in}}%
\pgfpathlineto{\pgfqpoint{1.287715in}{1.462048in}}%
\pgfpathlineto{\pgfqpoint{1.322657in}{1.433908in}}%
\pgfpathlineto{\pgfqpoint{1.353533in}{1.405830in}}%
\pgfpathlineto{\pgfqpoint{1.380863in}{1.378486in}}%
\pgfpathlineto{\pgfqpoint{1.416328in}{1.340019in}}%
\pgfpathlineto{\pgfqpoint{1.464347in}{1.284661in}}%
\pgfpathlineto{\pgfqpoint{1.502906in}{1.240507in}}%
\pgfpathlineto{\pgfqpoint{1.529257in}{1.212666in}}%
\pgfpathlineto{\pgfqpoint{1.547185in}{1.195832in}}%
\pgfpathlineto{\pgfqpoint{1.563953in}{1.182263in}}%
\pgfpathlineto{\pgfqpoint{1.579858in}{1.171772in}}%
\pgfpathlineto{\pgfqpoint{1.595142in}{1.164117in}}%
\pgfpathlineto{\pgfqpoint{1.610003in}{1.158923in}}%
\pgfpathlineto{\pgfqpoint{1.624591in}{1.155592in}}%
\pgfpathlineto{\pgfqpoint{1.662711in}{1.149397in}}%
\pgfpathlineto{\pgfqpoint{1.705917in}{1.159366in}}%
\pgfpathlineto{\pgfqpoint{1.720800in}{1.164279in}}%
\pgfpathlineto{\pgfqpoint{1.736074in}{1.171284in}}%
\pgfpathlineto{\pgfqpoint{1.751916in}{1.180940in}}%
\pgfpathlineto{\pgfqpoint{1.768554in}{1.193611in}}%
\pgfpathlineto{\pgfqpoint{1.786273in}{1.209571in}}%
\pgfpathlineto{\pgfqpoint{1.805428in}{1.229076in}}%
\pgfpathlineto{\pgfqpoint{1.826475in}{1.252416in}}%
\pgfpathlineto{\pgfqpoint{1.867470in}{1.300697in}}%
\pgfpathlineto{\pgfqpoint{1.907988in}{1.347734in}}%
\pgfpathlineto{\pgfqpoint{1.931752in}{1.372897in}}%
\pgfpathlineto{\pgfqpoint{1.958274in}{1.397554in}}%
\pgfpathlineto{\pgfqpoint{1.972629in}{1.409289in}}%
\pgfpathlineto{\pgfqpoint{1.987726in}{1.420548in}}%
\pgfpathlineto{\pgfqpoint{2.020186in}{1.441943in}}%
\pgfpathlineto{\pgfqpoint{2.055789in}{1.462771in}}%
\pgfpathlineto{\pgfqpoint{2.094832in}{1.484007in}}%
\pgfpathlineto{\pgfqpoint{2.137780in}{1.506187in}}%
\pgfpathlineto{\pgfqpoint{2.185247in}{1.529530in}}%
\pgfpathlineto{\pgfqpoint{2.237989in}{1.554114in}}%
\pgfpathlineto{\pgfqpoint{2.296914in}{1.579958in}}%
\pgfpathlineto{\pgfqpoint{2.363116in}{1.607071in}}%
\pgfpathlineto{\pgfqpoint{2.437909in}{1.635461in}}%
\pgfpathlineto{\pgfqpoint{2.522888in}{1.665147in}}%
\pgfpathlineto{\pgfqpoint{2.619997in}{1.696152in}}%
\pgfpathlineto{\pgfqpoint{2.673826in}{1.712161in}}%
\pgfpathlineto{\pgfqpoint{2.731623in}{1.728513in}}%
\pgfpathlineto{\pgfqpoint{2.793774in}{1.745214in}}%
\pgfpathlineto{\pgfqpoint{2.860713in}{1.762270in}}%
\pgfpathlineto{\pgfqpoint{2.932920in}{1.779687in}}%
\pgfpathlineto{\pgfqpoint{3.010936in}{1.797471in}}%
\pgfpathlineto{\pgfqpoint{3.095366in}{1.815629in}}%
\pgfpathlineto{\pgfqpoint{3.186887in}{1.834166in}}%
\pgfpathlineto{\pgfqpoint{3.286262in}{1.853090in}}%
\pgfpathlineto{\pgfqpoint{3.286262in}{1.853090in}}%
\pgfusepath{stroke}%
\end{pgfscope}%
\begin{pgfscope}%
\pgfpathrectangle{\pgfqpoint{0.484000in}{0.352669in}}{\pgfqpoint{3.330000in}{1.619031in}}%
\pgfusepath{clip}%
\pgfsetrectcap%
\pgfsetroundjoin%
\pgfsetlinewidth{0.803000pt}%
\definecolor{currentstroke}{rgb}{0.250980,0.400000,0.600000}%
\pgfsetstrokecolor{currentstroke}%
\pgfsetdash{}{0pt}%
\pgfpathmoveto{\pgfqpoint{0.635363in}{0.426291in}}%
\pgfpathlineto{\pgfqpoint{0.636441in}{0.950306in}}%
\pgfpathlineto{\pgfqpoint{0.647218in}{1.474322in}}%
\pgfpathlineto{\pgfqpoint{0.686605in}{1.725865in}}%
\pgfpathlineto{\pgfqpoint{0.744327in}{1.714367in}}%
\pgfpathlineto{\pgfqpoint{0.798634in}{1.697953in}}%
\pgfpathlineto{\pgfqpoint{0.849169in}{1.681610in}}%
\pgfpathlineto{\pgfqpoint{0.896225in}{1.665499in}}%
\pgfpathlineto{\pgfqpoint{0.940087in}{1.649624in}}%
\pgfpathlineto{\pgfqpoint{0.981014in}{1.633984in}}%
\pgfpathlineto{\pgfqpoint{1.054985in}{1.603401in}}%
\pgfpathlineto{\pgfqpoint{1.119781in}{1.573743in}}%
\pgfpathlineto{\pgfqpoint{1.176772in}{1.545008in}}%
\pgfpathlineto{\pgfqpoint{1.227102in}{1.517190in}}%
\pgfpathlineto{\pgfqpoint{1.271724in}{1.490251in}}%
\pgfpathlineto{\pgfqpoint{1.311432in}{1.464076in}}%
\pgfpathlineto{\pgfqpoint{1.346866in}{1.438375in}}%
\pgfpathlineto{\pgfqpoint{1.378514in}{1.412560in}}%
\pgfpathlineto{\pgfqpoint{1.393030in}{1.399345in}}%
\pgfpathlineto{\pgfqpoint{1.406718in}{1.385784in}}%
\pgfpathlineto{\pgfqpoint{1.419602in}{1.371819in}}%
\pgfpathlineto{\pgfqpoint{1.431709in}{1.357466in}}%
\pgfpathlineto{\pgfqpoint{1.453719in}{1.328053in}}%
\pgfpathlineto{\pgfqpoint{1.473062in}{1.298847in}}%
\pgfpathlineto{\pgfqpoint{1.497919in}{1.257952in}}%
\pgfpathlineto{\pgfqpoint{1.547910in}{1.173920in}}%
\pgfpathlineto{\pgfqpoint{1.562849in}{1.151709in}}%
\pgfpathlineto{\pgfqpoint{1.576543in}{1.134376in}}%
\pgfpathlineto{\pgfqpoint{1.585169in}{1.125575in}}%
\pgfpathlineto{\pgfqpoint{1.593509in}{1.119083in}}%
\pgfpathlineto{\pgfqpoint{1.601659in}{1.115088in}}%
\pgfpathlineto{\pgfqpoint{1.605692in}{1.114134in}}%
\pgfpathlineto{\pgfqpoint{1.609716in}{1.113977in}}%
\pgfpathlineto{\pgfqpoint{1.613745in}{1.114752in}}%
\pgfpathlineto{\pgfqpoint{1.617799in}{1.116683in}}%
\pgfpathlineto{\pgfqpoint{1.621901in}{1.120167in}}%
\pgfpathlineto{\pgfqpoint{1.626087in}{1.125963in}}%
\pgfpathlineto{\pgfqpoint{1.630420in}{1.135713in}}%
\pgfpathlineto{\pgfqpoint{1.635030in}{1.153995in}}%
\pgfpathlineto{\pgfqpoint{1.640008in}{1.170645in}}%
\pgfpathlineto{\pgfqpoint{1.644679in}{1.122883in}}%
\pgfpathlineto{\pgfqpoint{1.648862in}{1.122883in}}%
\pgfpathlineto{\pgfqpoint{1.656530in}{1.084919in}}%
\pgfpathlineto{\pgfqpoint{1.663269in}{1.063468in}}%
\pgfpathlineto{\pgfqpoint{1.669510in}{1.049460in}}%
\pgfpathlineto{\pgfqpoint{1.675434in}{1.039608in}}%
\pgfpathlineto{\pgfqpoint{1.681144in}{1.032575in}}%
\pgfpathlineto{\pgfqpoint{1.686706in}{1.027646in}}%
\pgfpathlineto{\pgfqpoint{1.692170in}{1.024379in}}%
\pgfpathlineto{\pgfqpoint{1.697572in}{1.022477in}}%
\pgfpathlineto{\pgfqpoint{1.702943in}{1.021729in}}%
\pgfpathlineto{\pgfqpoint{1.710996in}{1.022432in}}%
\pgfpathlineto{\pgfqpoint{1.719107in}{1.024990in}}%
\pgfpathlineto{\pgfqpoint{1.727340in}{1.029128in}}%
\pgfpathlineto{\pgfqpoint{1.738601in}{1.036761in}}%
\pgfpathlineto{\pgfqpoint{1.750302in}{1.046497in}}%
\pgfpathlineto{\pgfqpoint{1.765733in}{1.061210in}}%
\pgfpathlineto{\pgfqpoint{1.807720in}{1.102971in}}%
\pgfpathlineto{\pgfqpoint{1.815489in}{1.108930in}}%
\pgfpathlineto{\pgfqpoint{1.823439in}{1.113424in}}%
\pgfpathlineto{\pgfqpoint{1.831511in}{1.115755in}}%
\pgfpathlineto{\pgfqpoint{1.839621in}{1.115523in}}%
\pgfpathlineto{\pgfqpoint{1.847674in}{1.112599in}}%
\pgfpathlineto{\pgfqpoint{1.855586in}{1.107494in}}%
\pgfpathlineto{\pgfqpoint{1.867064in}{1.097360in}}%
\pgfpathlineto{\pgfqpoint{1.885010in}{1.078560in}}%
\pgfpathlineto{\pgfqpoint{1.907854in}{1.054840in}}%
\pgfpathlineto{\pgfqpoint{1.919960in}{1.043959in}}%
\pgfpathlineto{\pgfqpoint{1.931561in}{1.035462in}}%
\pgfpathlineto{\pgfqpoint{1.940017in}{1.030875in}}%
\pgfpathlineto{\pgfqpoint{1.948335in}{1.028047in}}%
\pgfpathlineto{\pgfqpoint{1.956587in}{1.027266in}}%
\pgfpathlineto{\pgfqpoint{1.962091in}{1.028074in}}%
\pgfpathlineto{\pgfqpoint{1.967630in}{1.030129in}}%
\pgfpathlineto{\pgfqpoint{1.973236in}{1.033642in}}%
\pgfpathlineto{\pgfqpoint{1.978951in}{1.038903in}}%
\pgfpathlineto{\pgfqpoint{1.984827in}{1.046322in}}%
\pgfpathlineto{\pgfqpoint{1.990933in}{1.056513in}}%
\pgfpathlineto{\pgfqpoint{1.997371in}{1.070485in}}%
\pgfpathlineto{\pgfqpoint{2.004299in}{1.090118in}}%
\pgfpathlineto{\pgfqpoint{2.012011in}{1.119859in}}%
\pgfpathlineto{\pgfqpoint{2.016353in}{1.142401in}}%
\pgfpathlineto{\pgfqpoint{2.021288in}{1.177292in}}%
\pgfpathlineto{\pgfqpoint{2.026671in}{1.183189in}}%
\pgfpathlineto{\pgfqpoint{2.032123in}{1.183189in}}%
\pgfpathlineto{\pgfqpoint{2.037365in}{1.165007in}}%
\pgfpathlineto{\pgfqpoint{2.042302in}{1.156083in}}%
\pgfpathlineto{\pgfqpoint{2.047096in}{1.151698in}}%
\pgfpathlineto{\pgfqpoint{2.051820in}{1.149408in}}%
\pgfpathlineto{\pgfqpoint{2.056510in}{1.148438in}}%
\pgfpathlineto{\pgfqpoint{2.061191in}{1.148533in}}%
\pgfpathlineto{\pgfqpoint{2.065883in}{1.149597in}}%
\pgfpathlineto{\pgfqpoint{2.070608in}{1.151585in}}%
\pgfpathlineto{\pgfqpoint{2.080229in}{1.158202in}}%
\pgfpathlineto{\pgfqpoint{2.090207in}{1.168140in}}%
\pgfpathlineto{\pgfqpoint{2.100704in}{1.181218in}}%
\pgfpathlineto{\pgfqpoint{2.111898in}{1.197338in}}%
\pgfpathlineto{\pgfqpoint{2.130457in}{1.227140in}}%
\pgfpathlineto{\pgfqpoint{2.159896in}{1.276099in}}%
\pgfpathlineto{\pgfqpoint{2.177272in}{1.301850in}}%
\pgfpathlineto{\pgfqpoint{2.186706in}{1.313939in}}%
\pgfpathlineto{\pgfqpoint{2.196635in}{1.325173in}}%
\pgfpathlineto{\pgfqpoint{2.207048in}{1.335581in}}%
\pgfpathlineto{\pgfqpoint{2.229290in}{1.354905in}}%
\pgfpathlineto{\pgfqpoint{2.266373in}{1.383949in}}%
\pgfpathlineto{\pgfqpoint{2.324290in}{1.427141in}}%
\pgfpathlineto{\pgfqpoint{2.376042in}{1.463400in}}%
\pgfpathlineto{\pgfqpoint{2.415604in}{1.489354in}}%
\pgfpathlineto{\pgfqpoint{2.460086in}{1.516702in}}%
\pgfpathlineto{\pgfqpoint{2.510405in}{1.545440in}}%
\pgfpathlineto{\pgfqpoint{2.567680in}{1.575587in}}%
\pgfpathlineto{\pgfqpoint{2.633281in}{1.607175in}}%
\pgfpathlineto{\pgfqpoint{2.669719in}{1.623526in}}%
\pgfpathlineto{\pgfqpoint{2.708905in}{1.640258in}}%
\pgfpathlineto{\pgfqpoint{2.751118in}{1.657381in}}%
\pgfpathlineto{\pgfqpoint{2.796672in}{1.674905in}}%
\pgfpathlineto{\pgfqpoint{2.845919in}{1.692841in}}%
\pgfpathlineto{\pgfqpoint{2.899255in}{1.711201in}}%
\pgfpathlineto{\pgfqpoint{2.957131in}{1.729999in}}%
\pgfpathlineto{\pgfqpoint{3.020055in}{1.749249in}}%
\pgfpathlineto{\pgfqpoint{3.088607in}{1.768967in}}%
\pgfpathlineto{\pgfqpoint{3.163446in}{1.789169in}}%
\pgfpathlineto{\pgfqpoint{3.245327in}{1.809873in}}%
\pgfpathlineto{\pgfqpoint{3.335114in}{1.831098in}}%
\pgfpathlineto{\pgfqpoint{3.433800in}{1.852864in}}%
\pgfpathlineto{\pgfqpoint{3.542533in}{1.875193in}}%
\pgfpathlineto{\pgfqpoint{3.662636in}{1.898108in}}%
\pgfpathlineto{\pgfqpoint{3.662636in}{1.898108in}}%
\pgfusepath{stroke}%
\end{pgfscope}%
\begin{pgfscope}%
\pgfpathrectangle{\pgfqpoint{0.484000in}{0.352669in}}{\pgfqpoint{3.330000in}{1.619031in}}%
\pgfusepath{clip}%
\pgfsetrectcap%
\pgfsetroundjoin%
\pgfsetlinewidth{0.803000pt}%
\definecolor{currentstroke}{rgb}{0.619608,0.188235,0.172549}%
\pgfsetstrokecolor{currentstroke}%
\pgfsetdash{}{0pt}%
\pgfpathmoveto{\pgfqpoint{0.635363in}{0.426299in}}%
\pgfpathlineto{\pgfqpoint{0.636441in}{0.950315in}}%
\pgfpathlineto{\pgfqpoint{0.647219in}{1.474330in}}%
\pgfpathlineto{\pgfqpoint{0.688121in}{1.737227in}}%
\pgfpathlineto{\pgfqpoint{0.748700in}{1.724975in}}%
\pgfpathlineto{\pgfqpoint{0.805468in}{1.707487in}}%
\pgfpathlineto{\pgfqpoint{0.858046in}{1.690073in}}%
\pgfpathlineto{\pgfqpoint{0.906777in}{1.672914in}}%
\pgfpathlineto{\pgfqpoint{0.951993in}{1.656011in}}%
\pgfpathlineto{\pgfqpoint{0.993996in}{1.639363in}}%
\pgfpathlineto{\pgfqpoint{1.033056in}{1.622966in}}%
\pgfpathlineto{\pgfqpoint{1.069420in}{1.606819in}}%
\pgfpathlineto{\pgfqpoint{1.103311in}{1.590922in}}%
\pgfpathlineto{\pgfqpoint{1.134932in}{1.575273in}}%
\pgfpathlineto{\pgfqpoint{1.164468in}{1.559876in}}%
\pgfpathlineto{\pgfqpoint{1.192087in}{1.544730in}}%
\pgfpathlineto{\pgfqpoint{1.217941in}{1.529841in}}%
\pgfpathlineto{\pgfqpoint{1.242171in}{1.515211in}}%
\pgfpathlineto{\pgfqpoint{1.264905in}{1.500847in}}%
\pgfpathlineto{\pgfqpoint{1.286261in}{1.486754in}}%
\pgfpathlineto{\pgfqpoint{1.306347in}{1.472941in}}%
\pgfpathlineto{\pgfqpoint{1.325262in}{1.459415in}}%
\pgfpathlineto{\pgfqpoint{1.343096in}{1.446184in}}%
\pgfpathlineto{\pgfqpoint{1.359934in}{1.433254in}}%
\pgfpathlineto{\pgfqpoint{1.375852in}{1.420626in}}%
\pgfpathlineto{\pgfqpoint{1.390920in}{1.408293in}}%
\pgfpathlineto{\pgfqpoint{1.405201in}{1.396230in}}%
\pgfpathlineto{\pgfqpoint{1.418752in}{1.384383in}}%
\pgfpathlineto{\pgfqpoint{1.431619in}{1.372650in}}%
\pgfpathlineto{\pgfqpoint{1.443837in}{1.360852in}}%
\pgfpathlineto{\pgfqpoint{1.455429in}{1.348725in}}%
\pgfpathlineto{\pgfqpoint{1.466405in}{1.335938in}}%
\pgfpathlineto{\pgfqpoint{1.476760in}{1.322194in}}%
\pgfpathlineto{\pgfqpoint{1.486486in}{1.307385in}}%
\pgfpathlineto{\pgfqpoint{1.495582in}{1.291704in}}%
\pgfpathlineto{\pgfqpoint{1.504065in}{1.275570in}}%
\pgfpathlineto{\pgfqpoint{1.511967in}{1.259454in}}%
\pgfpathlineto{\pgfqpoint{1.519335in}{1.243730in}}%
\pgfpathlineto{\pgfqpoint{1.526221in}{1.228637in}}%
\pgfpathlineto{\pgfqpoint{1.532674in}{1.214296in}}%
\pgfpathlineto{\pgfqpoint{1.538745in}{1.200755in}}%
\pgfpathlineto{\pgfqpoint{1.544474in}{1.188013in}}%
\pgfpathlineto{\pgfqpoint{1.549900in}{1.176047in}}%
\pgfpathlineto{\pgfqpoint{1.555056in}{1.164818in}}%
\pgfpathlineto{\pgfqpoint{1.559972in}{1.154278in}}%
\pgfpathlineto{\pgfqpoint{1.564671in}{1.144369in}}%
\pgfpathlineto{\pgfqpoint{1.569176in}{1.135018in}}%
\pgfpathlineto{\pgfqpoint{1.573503in}{1.126132in}}%
\pgfpathlineto{\pgfqpoint{1.577668in}{1.117581in}}%
\pgfpathlineto{\pgfqpoint{1.581681in}{1.109183in}}%
\pgfpathlineto{\pgfqpoint{1.585547in}{1.100679in}}%
\pgfpathlineto{\pgfqpoint{1.589268in}{1.091720in}}%
\pgfpathlineto{\pgfqpoint{1.592838in}{1.081887in}}%
\pgfpathlineto{\pgfqpoint{1.596249in}{1.070795in}}%
\pgfpathlineto{\pgfqpoint{1.599487in}{1.058261in}}%
\pgfpathlineto{\pgfqpoint{1.602543in}{1.044439in}}%
\pgfpathlineto{\pgfqpoint{1.605413in}{1.029756in}}%
\pgfpathlineto{\pgfqpoint{1.608103in}{1.014721in}}%
\pgfpathlineto{\pgfqpoint{1.610621in}{0.999753in}}%
\pgfpathlineto{\pgfqpoint{1.612980in}{0.985119in}}%
\pgfpathlineto{\pgfqpoint{1.615194in}{0.970962in}}%
\pgfpathlineto{\pgfqpoint{1.617278in}{0.957331in}}%
\pgfpathlineto{\pgfqpoint{1.619242in}{0.944231in}}%
\pgfpathlineto{\pgfqpoint{1.621098in}{0.931634in}}%
\pgfpathlineto{\pgfqpoint{1.622857in}{0.919503in}}%
\pgfpathlineto{\pgfqpoint{1.624525in}{0.907802in}}%
\pgfpathlineto{\pgfqpoint{1.626112in}{0.896490in}}%
\pgfpathlineto{\pgfqpoint{1.627622in}{0.885515in}}%
\pgfpathlineto{\pgfqpoint{1.629062in}{0.874836in}}%
\pgfpathlineto{\pgfqpoint{1.630437in}{0.864417in}}%
\pgfpathlineto{\pgfqpoint{1.631751in}{0.854224in}}%
\pgfpathlineto{\pgfqpoint{1.633008in}{0.844227in}}%
\pgfpathlineto{\pgfqpoint{1.634212in}{0.834399in}}%
\pgfpathlineto{\pgfqpoint{1.635365in}{0.824716in}}%
\pgfpathlineto{\pgfqpoint{1.636470in}{0.815154in}}%
\pgfpathlineto{\pgfqpoint{1.637530in}{0.805694in}}%
\pgfusepath{stroke}%
\end{pgfscope}%
\begin{pgfscope}%
\pgfsetrectcap%
\pgfsetmiterjoin%
\pgfsetlinewidth{0.602250pt}%
\definecolor{currentstroke}{rgb}{0.000000,0.000000,0.000000}%
\pgfsetstrokecolor{currentstroke}%
\pgfsetdash{}{0pt}%
\pgfpathmoveto{\pgfqpoint{0.484000in}{0.352669in}}%
\pgfpathlineto{\pgfqpoint{0.484000in}{1.971700in}}%
\pgfusepath{stroke}%
\end{pgfscope}%
\begin{pgfscope}%
\pgfsetrectcap%
\pgfsetmiterjoin%
\pgfsetlinewidth{0.602250pt}%
\definecolor{currentstroke}{rgb}{0.000000,0.000000,0.000000}%
\pgfsetstrokecolor{currentstroke}%
\pgfsetdash{}{0pt}%
\pgfpathmoveto{\pgfqpoint{0.484000in}{0.352669in}}%
\pgfpathlineto{\pgfqpoint{3.814000in}{0.352669in}}%
\pgfusepath{stroke}%
\end{pgfscope}%
\begin{pgfscope}%
\pgfsetrectcap%
\pgfsetroundjoin%
\pgfsetlinewidth{0.803000pt}%
\definecolor{currentstroke}{rgb}{0.768627,0.556863,0.211765}%
\pgfsetstrokecolor{currentstroke}%
\pgfsetdash{}{0pt}%
\pgfpathmoveto{\pgfqpoint{3.168736in}{1.301947in}}%
\pgfpathlineto{\pgfqpoint{3.259014in}{1.301947in}}%
\pgfpathlineto{\pgfqpoint{3.349292in}{1.301947in}}%
\pgfusepath{stroke}%
\end{pgfscope}%
\begin{pgfscope}%
\definecolor{textcolor}{rgb}{0.000000,0.000000,0.000000}%
\pgfsetstrokecolor{textcolor}%
\pgfsetfillcolor{textcolor}%
\pgftext[x=3.421514in,y=1.270350in,left,base]{\color{textcolor}{\sffamily\fontsize{6.500000}{7.800000}\selectfont\catcode`\^=\active\def^{\ifmmode\sp\else\^{}\fi}\catcode`\%=\active\def%{\%}$b = 8M$}}%
\end{pgfscope}%
\begin{pgfscope}%
\pgfsetrectcap%
\pgfsetroundjoin%
\pgfsetlinewidth{0.803000pt}%
\definecolor{currentstroke}{rgb}{0.250980,0.400000,0.600000}%
\pgfsetstrokecolor{currentstroke}%
\pgfsetdash{}{0pt}%
\pgfpathmoveto{\pgfqpoint{3.168736in}{1.169440in}}%
\pgfpathlineto{\pgfqpoint{3.259014in}{1.169440in}}%
\pgfpathlineto{\pgfqpoint{3.349292in}{1.169440in}}%
\pgfusepath{stroke}%
\end{pgfscope}%
\begin{pgfscope}%
\definecolor{textcolor}{rgb}{0.000000,0.000000,0.000000}%
\pgfsetstrokecolor{textcolor}%
\pgfsetfillcolor{textcolor}%
\pgftext[x=3.421514in,y=1.137843in,left,base]{\color{textcolor}{\sffamily\fontsize{6.500000}{7.800000}\selectfont\catcode`\^=\active\def^{\ifmmode\sp\else\^{}\fi}\catcode`\%=\active\def%{\%}$b \approx b_c$}}%
\end{pgfscope}%
\begin{pgfscope}%
\pgfsetrectcap%
\pgfsetroundjoin%
\pgfsetlinewidth{0.803000pt}%
\definecolor{currentstroke}{rgb}{0.619608,0.188235,0.172549}%
\pgfsetstrokecolor{currentstroke}%
\pgfsetdash{}{0pt}%
\pgfpathmoveto{\pgfqpoint{3.168736in}{1.035805in}}%
\pgfpathlineto{\pgfqpoint{3.259014in}{1.035805in}}%
\pgfpathlineto{\pgfqpoint{3.349292in}{1.035805in}}%
\pgfusepath{stroke}%
\end{pgfscope}%
\begin{pgfscope}%
\definecolor{textcolor}{rgb}{0.000000,0.000000,0.000000}%
\pgfsetstrokecolor{textcolor}%
\pgfsetfillcolor{textcolor}%
\pgftext[x=3.421514in,y=1.004207in,left,base]{\color{textcolor}{\sffamily\fontsize{6.500000}{7.800000}\selectfont\catcode`\^=\active\def^{\ifmmode\sp\else\^{}\fi}\catcode`\%=\active\def%{\%}$b = 4M$}}%
\end{pgfscope}%
\end{pgfpicture}%
\makeatother%
\endgroup%

  \caption[Adaptive step-size profiles near the photon sphere]{%
    Adaptive step size $h$ as a function of affine parameter
    $\lambda$ for three null geodesics in the Schwarzschild
    geometry ($M = 1$), integrated with RK45 at tolerances
    $\epsilon_a = \epsilon_r = 10^{-10}$.  Gold: a deflected ray
    ($b = 8M$) showing a brief step-size reduction near closest
    approach at $r \approx 6.7M$, followed by smooth recovery.
    Blue: a near-critical ray ($b \approx b_c$) whose step size
    drops by more than an order of magnitude and remains
    suppressed from $\lambda \approx 50$ to $\lambda \approx 75$
    as the ray winds near the photon sphere at $r \approx 3M$;
    the oscillations reflect per-orbit step-size variation.
    Crimson: a plunge ray ($b = 4M < b_c$) whose step size falls
    as it crosses the horizon, with no recovery.}
  \label{fig:step-size-profile}
\end{figure}

\paragraph{The integration pipeline.}

The function \texttt{integrateGeodesic} in
\codemodule{Geodesic.Solver} encapsulates the full ray-tracing
pipeline for a single pixel.  Its type signature separates the
concerns cleanly:
%
\begin{lstlisting}[language=Haskell]
integrateGeodesic
  :: (forall a. Floating a => [a] -> [a])
  -> (Vec4 'Contravariant -> Sym4 'Contravariant)
  -> SolverConfig
  -> TerminationConfig
  -> HamiltonianState
  -> GeodesicResult HamiltonianState
\end{lstlisting}
%
The first two arguments define the spacetime: the polymorphic metric
function (enabling AD) and the inverse metric lookup.  The solver and
termination configurations specify tolerances and stopping criteria
respectively.  The initial state $(x^\mu_0, p_{\mu,0},
\lambda_0 = 0)$ is set from the camera model and the null constraint
as described in \cref{sec:geo-null}.
\implnote{The \texttt{SolverConfig} record holds the seven
parameters governing step-size adaptation: initial step, absolute and
relative tolerances, maximum step count, safety factor, and the
minimum and maximum scaling bounds.}

The return type \texttt{GeodesicResult} pairs the final state with a
\texttt{TerminationReason} --- horizon crossing, escape to the
celestial sphere, accretion-disc intersection, or step-count
exhaustion (\cref{sec:geo-termination}).  This single algebraic data
type feeds directly into the rendering pipeline, where the
termination reason determines whether the pixel receives a disc
colour, a sky-map lookup, or a shadow fill.
\xrefnote{The rendering pipeline's use of \texttt{TerminationReason}
to assign pixel colours is developed in the ray-tracing chapter
(\cref{ch:ray-tracing}).}
